\NormalFont\ShortTitle{Atos}
{\MT Atos dos Apóstolos

\par }\ChapOne{1}{\PP \VerseOne{1}Eu fiz o primeiro livro, ó Teófilo, sobre todas as coisas que Jesus começou, tanto a fazer como a ensinar;
\VS{2}Até o dia em que ele foi recebido acima, depois de pelo Espírito Santo ter dado mandamentos aos apóstolos que tinha escolhido;
\VS{3}Aos quais também, depois de ter sofrido, apresentou-se vivo com muitas evidências; sendo visto por eles durante quarenta dias, e falando
{\ADD{-lhes}} das coisas relativas ao reino de Deus.
\VS{4}E, reunindo-os, mandou-lhes que não saíssem de Jerusalém, mas que esperassem a promessa do Pai que (disse ele) de mim ouvistes.
\VS{5}Porque João batizou com água, mas vós sereis batizados com o Espírito Santo, não muitos dias depois destes.
\VS{6}Então aqueles que tinham se reunido lhe perguntaram, dizendo: Senhor, tu restaurarás neste tempo o Reino a Israel?
\VS{7}E ele lhes disse: Não pertence a vós saber os tempos ou estações que o Pai pôs em sua própria autoridade.
\VS{8}Mas vós recebereis poder do Espírito Santo, que virá sobre vós; e vós sereis minhas testemunhas, tanto em Jerusalém como em toda a Judeia, e Samaria, e até ao último
{\ADD{lugar}} da terra.
\VS{9}E tendo ele dito estas coisas, enquanto eles o viam, ele foi levantado acima, e uma nuvem o tirou dos olhos deles.
\VS{10}E enquanto eles estavam com os olhos fixos ao céu, depois dele ter ido, eis que dois homens de roupas brancas se puseram junto a eles;
\VS{11}Os quais também disseram: Homens galileus, por que estais olhando para o céu? Este Jesus, que foi tomado de vós acima ao céu, assim virá, da maneira como o vistes ir ao céu.
\VS{12}Então eles voltaram a Jerusalém do monte que se chama das Oliveiras, o qual está perto de Jerusalém
{\ADD{à distância}} de um caminho de sábado.
\VS{13}E ao entrarem, subiram ao cômodo superior, onde ficaram Pedro, Tiago, João, André, Filipe, Tomé, Bartolomeu, Mateus, Tiago
{\ADD{filho}} de Alfeu, Simão Zelote e Judas
{\ADD{irmão}} de Tiago.
\VS{14}Todos estes perseveravam concordando em orações, e petições, com as mulheres, com Maria a mãe de Jesus, e com os irmãos dele.
\VS{15}E em
{\ADD{algum d}} aqueles dias, havendo uma multidão reunida de cerca de cento e vinte pessoas, Pedro se levantou no meio dos discípulos e disse:
\VS{16}Homens irmãos, era necessário que se cumprisse a Escritura, que o Espírito Santo, por meio da boca de Davi, predisse quanto a Judas, que foi o guia daqueles que prenderam a Jesus.
\VS{17}Porque ele foi contado conosco, e obteve uma porção neste ministério.
\VS{18}Este pois, adquiriu um campo por meio do pagamento da maldade, e tendo caído de cabeça para baixo, partiu-se ao meio, e todos os seus órgãos internos caíram para fora.
\VS{19}E
{\ADD{isso}} foi conhecido por todos os que habitam em Jerusalém, de maneira que aquele campo se chama em sua própria língua Aceldama, isto é, campo de sangue.
\VS{20}Porque está escrito no livro dos Salmos: Sua habitação se faça deserta, e não haja quem nela habite; e outro tome seu trabalho de supervisão.
\VS{21}Portanto é necessário, que dos homens que nos acompanharam todo o tempo em que o Senhor Jesus entrava e saía conosco,
\VS{22}Começando desde o batismo de João, até o dia em que
{\ADD{diante}} de nós ele foi recebido acima, se faça um destes testemunha conosco de sua ressurreição.
\VS{23}E apresentaram dois: a José, chamado Barsabás, que tinha por sobrenome
{\ADD{o}} Justo; e a Matias.
\VS{24}E orando, disseram: Tu, Senhor, conhecedor dos corações de todos, mostra a qual destes dois tu tens escolhido.
\VS{25}Para que ele tome parte deste ministério e apostolado, do qual Judas se desviou para ir a seu próprio lugar.
\VS{26}E lançaram-lhes as sortes; e caiu a sorte sobre Matias. E ele
{\ADD{passou}} a ser contado junto com os onze apóstolos.

\par }\Chap{2}{\PP \VerseOne{1}E ao se cumprir o dia de Pentecostes, estavam todos concordando no mesmo lugar.
\VS{2}E de repente houve um ruído do céu, como de um vento forte
{\ADD{e}} violento, e encheu toda a casa onde eles estavam sentados.
\VS{3}E foram vistas por eles línguas repartidas como que de fogo, e se pôs sobre cada um deles.
\VS{4}E eles foram todos cheios do Espírito Santo, e começaram a falar em outras línguas, conforme o Espírito lhes dava a discursarem.
\VS{5}E havia judeus que estavam morando em Jerusalém, homens devotos, de toda nação abaixo do céu.
\VS{6}E acontecendo esta voz, ajuntou-se a multidão; e ela estava confusa, porque cada um os ouvia falar em sua própria língua.
\VS{7}E todos estavam admirados, e se maravilhavam, dizendo uns aos outros: Ora, estes que estão falando, não são todos eles galileus?
\VS{8}E como nós ouvimos cada um
{\ADD{deles}} em nossa própria língua, na qual nascemos?
\VS{9}Partos, Medos, Elamitas, os habitantes da Mesopotâmia, da Judeia, Capadócia, Ponto, Ásia,
\VS{10}Frígia, Panfília, Egito e regiões da Líbia perto de Cirene, e romanos estrangeiros, tanto judeus como prosélitos,
\VS{11}Cretenses e Árabes, os ouvimos em nossas próprias línguas eles falarem das grandezas de Deus.
\VS{12}E todos estavam admirados e confusos, dizendo uns aos outros: O que isso quer dizer?
\VS{13}E outros, ridicularizando, diziam: Eles estão cheios de vinho doce.
\VS{14}Mas Pedro, pondo-se de pé com os onze, levantou sua voz, e lhes falou: Homens judeus, e todos os que habitais em Jerusalém, seja isto conhecido, e ouvi minhas palavras:
\VS{15}Porque estes não estão bêbados, como vós pensais, sendo
{\ADD{ainda}} a terceira hora do dia.
\VS{16}Mas isto é o que foi dito por meio do profeta Joel:
\VS{17}E será nos últimos dias, diz Deus, que: Eu derramarei do meu Espírito sobre toda carne, e vossos filhos e filhas profetizarão, e vossos rapazes terão visões, e vossos velhos sonharão sonhos;
\VS{18}E também sobre meus servos e sobre minhas servas, naqueles dias eu derramarei do meu Espírito, e profetizarão.
\VS{19}E darei milagres acima no céu, e sinais abaixo na terra; sangue, fogo, e vapor de fumaça;
\VS{20}O sol se converterá em trevas, e a lua em sangue, antes que venha o grande e notório dia do Senhor.
\VS{21}E será que todo aquele que chamar ao nome do Senhor será salvo.
\VS{22}Homens israelitas, ouvi estas palavras: Jesus o nazareno, homem aprovado por Deus entre vós, com maravilhas, milagres e sinais, que Deus fez por meio dele no meio de vós, assim como vós mesmos também sabeis;
\VS{23}Este, sendo entregue pelo determinado conselho e conhecimento prévio de Deus, sendo tomando, pelas mãos de injustos
{\ADD{o}} crucificastes e matastes;
\VS{24}Ao qual Deus ressuscitou, tendo soltado as dores da morte; porque não era possível ele ser retido por ela;
\VS{25}Porque Davi diz sobre ele: Eu sempre via ao Senhor diante de mim, porque ele está à minha direita, para que eu não seja abalado.
\VS{26}Por isso meu coração está contente, e minha língua se alegra, e até mesmo minha carne repousará em esperança.
\VS{27}Pois tu não abandonarás minha alma no Xeol,
\FTNT{*}{{\FR 2:27 }
Xeol é o lugar dos mortos
} nem permitirás que o teu santo veja a degradação.
\FTNT{†}{{\FR 2:27 }
Ou: deterioração, putrefação. Também no v. 31
}
\VS{28}Tu tens me feito conhecer os caminhos da vida; tu me encherás de alegria com tua face.
\VS{29}Homens irmãos, é lícito eu vos dizer abertamente sobre o patriarca Davi, que morreu, e foi sepultado, e a sepultura dele está conosco até o dia de hoje.
\VS{30}Portanto, sendo ele profeta, e sabendo que Deus tinha lhe prometido com juramento que, da sua descendência
\FTNT{‡}{{\FR 2:30 }
da sua descendência
lit. do fruto de seus lombos, i.e., de um descendente biológico
} segundo a carne, levantaria ao Cristo para se sentar no seu trono;
\VS{31}Vendo
{\ADD{-o}} com antecedência, falou da ressurreição do Cristo, que a sua alma não foi abandonada no Xeol, nem a sua carne viu a degradação.
\VS{32}A este Jesus Deus ressuscitou; do qual todos nós somos testemunhas.
\VS{33}Portanto, tendo sido exaltado à direita de Deus, e recebido do Pai a promessa do Espírito Santo, derramou isto que agora estais vendo e ouvindo.
\VS{34}Porque Davi não subiu aos céus; mas sim, ele diz: Disse o Senhor a meu Senhor: Senta-te à minha direita,
\VS{35}Até que eu ponha teus inimigos por escabelo de teus pés.
\VS{36}Saiba então com certeza toda a casa de Israel, que Deus o fez Senhor e Cristo a este Jesus, que vós crucificastes.
\VS{37}E eles, ao ouvirem
{\ADD{estas coisas}} , foram afligidos como que perfurados de coração, e disseram a Pedro, e aos outros apóstolos: Que faremos, homens irmãos?
\VS{38}E Pedro lhes disse: Arrependei-vos, e batize-se cada um de vós no nome de Jesus Cristo, para perdão dos pecados; e vós recebereis o dom do Espírito Santo.
\VS{39}Porque a promessa é para vós, e para vossos filhos, e para todos que
{\ADD{ainda}} estão longe, a tantos quantos Deus, nosso Senhor, chamar.
\VS{40}E com muitas outras palavras ele dava testemunho, e exortava, dizendo: Salvai-vos desta geração perversa!
\VS{41}Então os que receberam a palavra dele de boa vontade foram batizados; e foram adicionados naquele dia quase três mil almas.
\VS{42}E eles perseveravam na doutrina dos apóstolos, na comunhão, no partir do pão, e nas orações.
\VS{43}E houve temor em toda alma; e muitos milagres e sinais foram feitos pelos apóstolos.
\VS{44}E todos os que criam estavam juntos, e tinham todas as coisas em comum.
\VS{45}E eles vendiam
{\ADD{suas}} propriedades e bens, e as repartiam com todos, conforme a necessidade que cada um tinha.
\VS{46}E perseverando a cada dia em concordância no Templo, e partindo o pão de casa em casa, comiam juntos com alegria e sinceridade de coração.
\VS{47}Louvando a Deus, e tendo graça,
{\ADD{sendo do agrado}} de todo o povo. E a cada dia o Senhor acrescentava à igreja aqueles que estavam sendo salvos.

\par }\Chap{3}{\PP \VerseOne{1}E Pedro e João estavam subindo juntos para o Templo à hora da oração (a nona
{\ADD{hora}} );
\VS{2}E um certo homem estava sendo trazido, que era aleijado desde o ventre de sua mãe, ao qual todo dia colocavam à porta do Templo, chamada
{\ADD{Porta}} Formosa, para pedir esmola aos que entravam no Templo.
\VS{3}O qual, ao ver Pedro e João perto de entrarem no Templo, ele
{\ADD{lhes}} pediu uma esmola.
\VS{4}E Pedro, olhando fixamente para ele, junto com João, disse: Olha para nós.
\VS{5}E
{\ADD{o aleijado}} ficou prestando atenção neles, esperando receber deles alguma coisa.
\VS{6}E Pedro disse: Prata e ouro eu não tenho; mas o que eu tenho, isso eu te dou: no nome de Jesus Cristo, o nazareno, levanta-te, e anda!
\VS{7}E, tomando-o pela mão direita, levantou
{\ADD{-o}} ; e logo os seus pés e tornozelos ficaram firmes.
\VS{8}E ele, saltando, pôs-se de pé, e andou, e entrou com eles no Templo, andando, e saltando, e louvando a Deus.
\VS{9}E todo o povo o viu andar, e louvar a Deus.
\VS{10}E eles o reconheceram, que este era o que se sentava
{\ADD{para pedir}} esmola perto da porta formosa do Templo; e ficaram cheios de surpresa e espanto, por causa do que tinha lhe acontecido.
\VS{11}E o aleijado que tinha sido curado, tendo se apegado a Pedro e a João, todo o povo correu maravilhado a eles ao pórtico, que se chama de Salomão.
\VS{12}Quando Pedro viu
{\ADD{isso}} , respondeu ao povo: “Homens israelitas, por que vos maravilhais disto? Ou por que vós olhais tão atentamente para nós, como se por nosso próprio poder ou devoção divina o tivéssemos feito andar?
\VS{13}O Deus de Abraão, e de Isaque, e de Jacó, o Deus de nossos pais, glorificou a seu filho Jesus, ao qual vós entregastes, e diante do rosto de Pilatos o negastes,
{\ADD{mesmo}} ele julgando que fosse solto.
\VS{14}Mas vós negastes ao santo e justo, e pedistes que um homem assassino fosse vos dado.
\VS{15}E vós matastes ao Príncipe da vida, ao qual Deus ressuscitou dos mortos, do que nós somos testemunhas.
\VS{16}E pela fé em seu nome, o nome dele deu firmeza a este, que vedes e conheceis; e a fé que é por meio dele deu a este perfeita saúde na presença de todos vós.
\VS{17}E agora, irmãos, eu sei que vós fizestes
{\ADD{isso}} por ignorância, assim como também vossos líderes.
\VS{18}Mas Deus cumpriu assim o que já antes pela boca de todos os seus profetas ele tinha anunciado, que o Cristo tinha de sofrer.
\VS{19}Arrependei-vos, pois, e convertei-vos, para que vosso pecados sejam apagados, quando vierem os tempos do refrigério da presença do Senhor.
\VS{20}E ele enviar a Jesus Cristo, que já vos foi pregado anteriormente.
\VS{21}Ao qual convém que o céu receba até os tempos da restauração de todas as coisas, que Deus falou pela boca de todos os seus santos profetas, desde o princípio.
\VS{22}Porque Moisés disse aos
{\ADD{nossos}} pais: O Senhor, vosso Deus, levantará dentre vossos irmãos um profeta como eu; a ele ouvireis em tudo o que ele vos falar.
\VS{23}E será que toda pessoa
\FTNT{*}{{\FR 3:23 }
Lit. alma
} que não ouvir este profeta será exterminada do povo.
\VS{24}E também todos os profetas, desde Samuel e os posteriores, todos os que falaram, também anunciaram com antecedência destes dias.
\VS{25}Vós sois os filhos dos profetas e do pacto que Deus estabeleceu com nossos pais, dizendo a Abraão: E em tua semente serão abençoadas todas as famílias da terra.
\VS{26}Deus, ao ressuscitar seu filho Jesus, primeiro o enviou a vós, para que
{\ADD{nisto}} vos abençoasse: afastando cada um de vós de vossas maldades”.

\par }\Chap{4}{\PP \VerseOne{1}E enquanto eles ainda estavam falando ao povo, vieram sobre eles os sacerdotes, e o oficial do Templo, e os saduceus,
\VS{2}Muito incomodados por eles ensinarem ao povo, e anunciarem no
{\ADD{nome}} de Jesus a ressurreição dos mortos.
\VS{3}E puseram as mãos sobre eles, e os puseram na prisão até o dia seguinte, porque já era tarde.
\VS{4}E muitos dos que ouviram a palavra, creram; e era o número dos homens de cerca de cinco mil.
\VS{5}E aconteceu no dia seguinte, que os chefes, e anciãos, e escribas, se reuniram em Jerusalém;
\VS{6}E Anás, o sumo sacerdote, e Caifás, e João, e Alexandre, e todos quantos havia da linhagem do governo sacerdotal.
\VS{7}E pondo-os no meio, perguntaram
{\ADD{-lhes}} : Por meio de que poder ou por qual nome vós fizestes isto?
\VS{8}Então Pedro, cheio do Espírito Santo, lhes disse: Chefes do povo, e anciões de Israel,
\VS{9}Se hoje somos interrogados quanto a uma boa ação
{\ADD{feita}} a um enfermo, pela qual este foi curado;
\VS{10}Seja conhecido a todos vós, e a todos o povo de Israel, que pelo nome de Jesus Cristo, o nazareno, aquele que vós crucificastes, ao qual Deus ressuscitou dos mortos, por ele este
{\ADD{homem}} está são diante de vós.
\VS{11}Este é a pedra que foi desprezada por vós, edificadores; a qual foi feita por cabeça de esquina.
\VS{12}E em nenhum outro há salvação; porque nenhum outro nome há abaixo do céu, dado entre os seres humanos, em quem devemos ser salvos.
\VS{13}E eles, ao verem a ousadia de Pedro, e de João; e informados que eles eram homens sem instrução e ordinários, maravilharam-se; e eles sabiam que eles tinham estado com Jesus.
\VS{14}E vendo estar com eles o homem que tinha sido curado, nada tinham a dizer contra
{\ADD{eles}} .
\VS{15}E mandando-os saírem do supremo conselho,
\FTNT{*}{{\FR 4:15 }
supremo conselho
lit. sinédrio – o mais importante conselho ou tribunal para os judeus
} discutiam entre si,
\VS{16}Dizendo: O que faremos a estes homens? Porque um sinal notório foi feito por eles, manifesto a todos os que habitam em Jerusalém, e não podemos negar.
\VS{17}Mas para que
{\ADD{esta notícia}} não seja ainda mais divulgada entre o povo, façamos sérias ameaças a eles, para que nunca mais falem a ninguém neste nome.
\VS{18}E chamando-os, ordenaram-lhes que nunca mais falassem nem ensinassem no nome de Jesus.
\VS{19}Mas, respondendo Pedro, e João, disseram-lhes: Julgai se é justo diante de Deus, ouvir a vós mais do que a Deus;
\VS{20}Porque nós não podemos deixar de falar daquilo que temos visto e ouvido.
\VS{21}Mas eles, tendo os ameaçado ainda mais, nada acharam
{\ADD{de motivo}} para os castigar, e os deixaram ir por causa do povo; porque todos glorificavam a Deus pelo que tinha acontecido.
\VS{22}Porque era de mais de quarenta anos o homem em quem este milagre de cura tinha sido feito.
\VS{23}E eles, tendo sido soltos, vieram aos seus
{\ADD{companheiros}} , e lhes contaram tudo o que os chefes dos sacerdotes e os anciãos tinham lhes dito.
\VS{24}E eles, ao ouvirem
{\ADD{isto}} , levantaram concordantes a voz a Deus, e disseram: Senhor, tu és o Deus que fizeste o céu, a terra, o mar, e todas as coisas que neles há;
\VS{25}Que pela boca de teu servo Davi disseste: Por que os gentios se irritam, e os povos gastam seus pensamentos em coisas vãs?
\VS{26}Os reis da terra se levantaram, e os príncipes se juntaram em um mesmo
{\ADD{propósito}} contra o Senhor, e contra o seu Ungido.
\FTNT{†}{{\FR 4:26 }
Ungido = equiv. Cristo
}
\VS{27}Porque verdadeiramente contra teu Santo Filho Jesus, ao qual tu ungiste, se ajuntaram, tanto Herodes, como Pôncio Pilatos, com os gentios e os povos de Israel.
\VS{28}Para fazerem tudo o que a tua mão e o teu conselho desde antes tinha determinado para acontecer.
\VS{29}E agora, Senhor, observa as ameaças deles, e dá a teus servos, que com toda ousadia falem tua palavra;
\VS{30}Estendendo tua mão para a cura, e que se façam sinais e milagres pelo nome de teu Santo filho Jesus.
\VS{31}E tendo orado, agitou-se o lugar em que eles estavam juntos, e foram todos cheios do Espírito Santo, e falavam a palavra de Deus com ousadia.
\VS{32}E a multidão dos que criam, era de um só oração e uma só alma; e ninguém dizia ser próprio coisa alguma de seus bens, mas todas as coisas lhes eram comuns.
\VS{33}E com grande poder os apóstolos davam testemunho da ressurreição do Senhor Jesus; e em todos eles havia grande graça.
\VS{34}Porque também nenhum necessitado havia entre eles; porque todos os que possuíam propriedades de terras, ou casas, vendendo
{\ADD{-as}} , traziam o valor das coisas vendidas, e
{\ADD{o}} depositavam junto aos pés dos apóstolos.
\VS{35}E a cada um se repartia segundo cada qual tinha necessidade.
\VS{36}E José, chamado pelos apóstolos pelo sobrenome de Barnabé (que traduzido é filho da consolação), levita, natural do Chipre,
\VS{37}Tendo ele uma propriedade de terra, vendeu
{\ADD{-a}} , e trouxe o valor, e
{\ADD{o}} depositou junto aos pés dos apóstolos.

\par }\Chap{5}{\PP \VerseOne{1}E um certo homem, de nome Ananias, com Safira, sua mulher, vendeu uma propriedade de terra.
\VS{2}E escondeu
{\ADD{parte}} do valor, sabendo também a mulher dele; e trazendo uma certa parte, depositou
{\ADD{-a}} junto aos pés dos apóstolos.
\VS{3}E Pedro disse: Ananias, por que Satanás encheu teu coração, para que mentisses ao Espírito Santo, e escondesses
{\ADD{parte}} do valor da propriedade?
\VS{4}{\ADD{Se}} mantivesses
{\ADD{tua propriedade}} , ela não seria mantida contigo? E, tendo sido vendida,
{\ADD{o dinheiro da venda}} não estava em teu poder? Por que decidiste
{\ADD{isto}} em teu coração? Não mentiste aos seres humanos, mas sim a Deus.
\VS{5}E Ananias, ao ouvir estas palavras, caiu e deixou de respirar. E veio um grande temor sobre todos os que ouviram isso.
\VS{6}E os rapazes, tendo se levantado, levaram-no fora, e
{\ADD{o}} sepultaram.
\VS{7}E passando o intervalo de cerca de três horas, entrou também sua mulher, não sabendo o que tinha acontecido.
\VS{8}E Pedro disse a ela: Dize-me, vendestes por
{\ADD{aquele}} tanto aquela propriedade?E ela disse: Sim, por
{\ADD{aquele}} tanto.
\VS{9}E Pedro lhe disse: Por que vós fizestes acordo para tentar ao Espírito do Senhor? Eis que
{\ADD{estão}} à porta os pés daqueles que sepultaram a teu marido, e eles também
{\ADD{te}} levarão.
\VS{10}E imediatamente ela caiu junto aos pés deles, e deixou de respirar. E os rapazes, ao entrarem, encontraram-na morta; e levando-
{\ADD{a}} fora, sepultaram-na junto ao marido dela.
\VS{11}E veio um grande temor em toda a igreja, e em todos que ouviram estas coisas.
\VS{12}E pelas mãos dos apóstolos foram feitos muitos sinais e milagres entre o povo. E estavam todos em concordância no pórtico de Salomão.
\VS{13}Mas dos outros, ninguém ousava se juntar a eles; porém o povo os estimava grandemente.
\VS{14}E cada vez mais os que criam no Senhor se aumentavam, multidões tanto de homens como de mulheres.
\VS{15}De maneira que traziam os enfermos às ruas, e os botavam em camas e macas, para que, vindo Pedro, pelo menos a sombra dele cobrisse a alguns deles.
\VS{16}E também das cidades vizinhas vinha uma multidão a Jerusalém, trazendo enfermos, e atormentados por espíritos imundos, os quais todos eram curados.
\VS{17}E levantando-se o sumo sacerdote, e todos os que estavam com ele (que eram do grupo sectário dos saduceus), eles se encheram de inveja.
\VS{18}E puseram suas mãos nos apóstolos, e os colocaram na prisão pública.
\VS{19}Mas um anjo do Senhor, durante a noite, abriu as portas da prisão; e levando-os para fora, disse:
\VS{20}Ide; ficai em pé, e falai no Templo ao povo todas as palavras desta vida.
\VS{21}E eles, ouvindo
{\ADD{isto}} , entraram de manhã cedo no templo, e ensinavam. Mas vindo o sumo sacerdote, e os que estavam com ele, chamaram ao supremo conselho,
\FTNT{*}{{\FR 5:21 }
supremo conselho
lit. sinédrio – o mais importante conselho ou tribunal para os judeus
} e a todos os anciãos dos filhos de Israel, e mandaram à prisão, para que os trouxessem.
\VS{22}Mas quando os oficiais vieram, não os acharam na prisão; e voltando, anunciaram,
\VS{23}Dizendo: Nós achamos a prisão fechada, em toda segurança, e com os guardas que estavam fora junto às portas; mas quando
{\ADD{as}} abrimos, a ninguém achamos dentro.
\VS{24}Quando o
{\ADD{sumo}} sacerdote, o chefe da guarda do Templo, e os chefes dos sacerdotes ouviram estas palavras, eles duvidaram deles quanto o que aquilo viria a ser.
\VS{25}E vindo alguém, anunciou-lhes, dizendo: Eis que os homens que vós pusestes na prisão estão no Templo, e ensinam ao povo.
\VS{26}Então foi o chefe da guarda do Templo com os oficiais, e os trouxe,
{\ADD{mas}} não com violência, porque temiam ao povo, para que não fossem apedrejados.
\VS{27}E quando os trouxeram, apresentaram-nos ao supremo conselho. E o sumo sacerdote perguntou a eles, dizendo:
\VS{28}Não vos ordenamos expressamente para não ensinardes
{\ADD{mais}} neste nome? E eis que vós enchestes a Jerusalém com vossa doutrina, e quereis trazer sobre nós o sangue deste homem!
\VS{29}E Pedro, respondendo com os apóstolos, disseram: Maior obrigação é obedecer a Deus do que às pessoas.
\VS{30}O Deus de nossos Pais ressuscitou a Jesus, ao qual vós matastes, pendurando
{\ADD{-o}} no madeiro.
\VS{31}A este Deus exaltou com sua
{\ADD{mão}} direita
{\ADD{por}} Príncipe e Salvador, para dar a Israel arrependimento e perdão de pecados.
\VS{32}E nós somos testemunhas dele quanto a estas palavras, e também o Espírito Santo, o qual Deus tem dado a aqueles que lhe obedecem.
\VS{33}E eles, ouvindo
{\ADD{isso}} , enfureceram-se grandemente, e planejaram matá-los.
\VS{34}Mas, levantando-se no supremo conselho um certo fariseu, de nome Gamaliel, instrutor da Lei, bem honrado por todo o povo, ele mandou levarem aos apóstolos para fora por um pouco
{\ADD{de tempo}} .
\VS{35}E lhes disse: Homens israelitas, olhai por vós mesmos, quanto ao que haveis de fazer a estes homens;
\VS{36}Porque antes destes dias se levantou Teudas, dizendo ser alguém; ao qual se ajuntaram cerca de quatrocentos homens; ao qual foi morto, e todos os que acreditavam
{\ADD{nele}} foram dispersos, e reduzidos a nada.
\VS{37}Depois deste se levantou Judas, o galileu, nos dias do censo; e perverteu muito do povo atrás dele; e este também pereceu, e todos os que acreditavam nele foram dispersos.
\VS{38}E agora, eu vos digo, afastai-vos destes homens, e deixai-os; porque se este conselho ou esta obra for humana, ela se desfará.
\VS{39}Mas se é de Deus, vós não a podereis desfazer; para que não venhais a ser achados de também lutardes contra Deus.
\VS{40}E concordaram com ele. E chamando aos apóstolos, tendo
{\ADD{os}} açoitado, mandaram
{\ADD{-lhes}} que não
{\ADD{mais}} falassem no nome de Jesus; e os deixaram ir.
\VS{41}Então eles saíram da presença do supremo conselho, alegres por terem sido considerados dignos de serem humilhados por causa do nome dele.
\VS{42}E todos os dias no Templo, e pelas casas, não paravam de ensinar e de anunciar a Jesus Cristo.

\par }\Chap{6}{\PP \VerseOne{1}Naqueles dias, ao se multiplicar
{\ADD{o número}} de discípulos, houve uma murmuração dos gregos contra os hebreus, de que suas viúvas estavam sendo desprezadas no serviço diário
{\ADD{de entrega de comida}} .
\VS{2}E os doze, chamando à multidão dos discípulos, disseram: Não é bom que nós deixemos a palavra de Deus para servirmos às mesas.
\VS{3}Portanto, irmãos, buscai sete homens dentre vós, de quem haja
{\ADD{bom}} testemunho, cheios do Espírito Santo e de sabedoria, aos quais constituamos sobre esta importante tarefa.
\VS{4}Nós, porém, perseveraremos na oração e no serviço da palavra.
\VS{5}E esta palavra foi do agrado diante de toda a multidão, e escolheram a Estêvão, homem cheio de fé e do Espírito Santo, e Filipe, Prócoro, Nicanor, Timom, Parmenas e a Nicolão, o prosélito de Antioquia.
\VS{6}Aos quais se apresentaram diante dos apóstolos; e eles, orando, puseram as mãos sobre eles.
\VS{7}E a palavra de Deus crescia, e o número dos discípulos se multiplicava muito em Jerusalém; e grande multidão dos sacerdotes obedecia à fé.
\VS{8}E Estêvão, cheio de fé e poder, fazia milagres e grandes sinais entre o povo.
\VS{9}E levantaram-se alguns da sinagoga,
{\ADD{que era}} chamada
{\ADD{sinagoga}} dos libertos, cireneus, e alexandrinos, e dos
{\ADD{que eram}} da Cilícia, e da Ásia, e discutiam contra Estêvão.
\VS{10}E eles não podiam resistir à sabedoria e ao Espírito com que ele falava.
\VS{11}Então eles subornaram a uns homens,
{\ADD{para que}} dissessem: Nós o ouvimos falando palavras blasfemas contra Moisés e
{\ADD{contra}} Deus.
\VS{12}E incitaram ao povo, aos anciãos e aos escribas; e vieram sobre
{\ADD{Estêvão}} , e o detiveram, e o levaram ao supremo conselho.
\FTNT{*}{{\FR 6:12 }
supremo conselho
lit. sinédrio – o mais importante conselho ou tribunal para os judeus
}
\VS{13}E apresentaram falsas testemunhas, que diziam: Este homem não para de falar palavras blasfemas contra este santo lugar, e
{\ADD{contra}} a Lei.
\VS{14}Porque nós o ouvimos dizer que este Jesus Nazareno vai destruir este lugar, e mudar os costumes que Moisés nos entregou.
\VS{15}Então todos os que estavam sentados no supremo conselho, observando-o com atenção, viram o rosto dele como de um anjo.

\par }\Chap{7}{\PP \VerseOne{1}E disse o chefe dos sacerdotes: Por acaso é isto assim?
\VS{2}E ele disse: Homens irmãos e pais, ouvi: o Deus da glória apareceu a nosso pai Abraão, estando ele na Mesopotâmia, antes de habitar em Harã;
\VS{3}E disse-lhe: Sai de tua terra, e de tua parentela, e vem a terra que eu te mostrarei.
\VS{4}Então ele saiu da terra dos caldeus,
{\ADD{e}} habitou em Harã. E dali, depois que morreu seu pai, ele partiu para esta terra, em que agora vós habitais.
\VS{5}E
{\ADD{Deus}} não lhe deu herança nela, nem mesmo a pegada de um pé; mas prometeu que a daria a ele em propriedade, e a sua semente depois dele, não tendo ele filho
{\ADD{ainda}} .
\VS{6}E Deus falou assim: Tua semente será peregrina em terra alheia, e a escravizarão, e a maltratarão
{\ADD{por}} quatrocentos anos.
\VS{7}E à nação a quem eles servirem, eu a julgarei, (disse Deus). E depois disso eles sairão, e me servirão neste lugar.
\VS{8}E ele lhe deu o pacto da circuncisão; e assim gerou a Isaque, e o circuncidou ao oitavo dia; e Isaque
{\ADD{gerou}} a Jacó, e Jacó aos doze patriarcas.
\VS{9}E os patriarcas, tendo inveja
{\ADD{de}} José, venderam
{\ADD{-no}} ao Egito; mas Deus era com ele.
\VS{10}E o livrou de todas as suas aflições, e lhe deu graça e sabedoria diante de Faraó, rei do Egito; e o pôs
{\ADD{por}} governador sobre o Egito, e toda a sua casa.
\VS{11}E veio fome sobre toda a terra do Egito e de Canaã, e grande aflição; e nossos pais não achavam alimentos.
\VS{12}Mas Jacó, ao ouvir que havia cereal no Egito, ele enviou nossos pais a primeira vez.
\VS{13}E na segunda
{\ADD{vez}} ,José foi reconhecido pelos seus irmãos, e a família de José foi conhecia por Faraó.
\VS{14}E José mandou chamar a seu pai Jacó, e toda a sua parentela, setenta e cinco almas.
\VS{15}E Jacó desceu ao Egito, e morreu; ele, e nossos pais;
\VS{16}E foram levados a Siquém, e postos na sepultura que Abraão, por uma quantia em dinheiro, tinha comprado dos filhos de Emor,
{\ADD{pai}} de Siquém.
\VS{17}Mas quando chegou perto o tempo da promessa que Deus tinha prometido a Abraão, o povo cresceu e se multiplicou no Egito.
\VS{18}Até que se levantou outro rei, que não tinha conhecido a José.
\VS{19}Este, usando de astúcia para com nossa parentela, maltratou a nossos pais, até fazendo com que eles rejeitassem suas crianças, para que não sobrevivessem.
\VS{20}Naquele tempo nasceu Moisés, e ele era muito formoso para Deus, e ele foi criado por três meses na casa de seu pai.
\VS{21}E tendo sido abandonado, a filha de Faraó o tomou, e o criou para si como filho.
\VS{22}E Moisés foi instruído em toda a sabedoria dos egípcios; e era poderoso em palavras e ações.
\VS{23}E quando lhe foi completado o tempo de quarenta anos
{\ADD{de idade}} , veio ao seu coração
{\ADD{o desejo}} de visitar a seus irmãos, os filhos de Israel.
\VS{24}E vendo um
{\ADD{deles}} sofrendo injustamente, defendeu
{\ADD{-o}} , e vingou pelo que tinha sido oprimido, matando ao egípcio.
\VS{25}E ele pensava que seus irmãos tivessem entendido que Deus ia lhes dar liberdade por meio da mão dele; mas eles não entenderam.
\VS{26}E no dia seguinte, estando
{\ADD{uns deles}} lutando, ele foi visto por eles, e ordenou-lhes
{\ADD{fazerem}} as pazes, dizendo: Homens, vós sois irmãos, por que fazeis mal um ao outro?
\VS{27}Mas aquele que maltratava a seu próximo empurrou-o, dizendo: Quem te pôs por chefe e juiz sobre nós?
\VS{28}Queres tu
{\ADD{também}} matar a mim, assim como ontem mataste ao egípcio?
\VS{29}E com esta palavra Moisés fugiu, e foi peregrino na terra de Midiã, onde ele gerou dois filhos.
\VS{30}E completados quarenta anos, um anjo do Senhor lhe apareceu no deserto do monte Sinai, em uma sarça inflamada.
\VS{31}Moisés, ao ver
{\ADD{isso}} ,maravilhou-se da visão; e ao aproximar-se para ver, veio até ele a voz do Senhor,
\VS{32}{\ADD{Dizendo}} : Eu
{\ADD{sou}} o Deus de teus pais, o Deus de Abraão, o Deus de Isaque, e o Deus de Jacó. E Moisés, estando tremendo, não ousava olhar com atenção.
\VS{33}E o Senhor lhe disse: Descalça-te as sandálias de teus pés, porque o lugar em que estás é terra santa.
\VS{34}Eu tenho visto com atenção a aflição de meu povo que
{\ADD{está}} no Egito, e ouvi o gemido deles, e eu desci para livrá-los; então vem agora, eu te enviarei ao Egito.
\VS{35}A este Moisés, ao qual tinham negado, dizendo: Quem te pôs por chefe e juiz, A este Deus enviou por chefe e libertador, pela mão do anjo que lhe apareceu na sarça.
\VS{36}Este os levou para fora, fazendo milagres e sinais na terra do Egito, e no mar Vermelho, e no deserto,
{\ADD{por}} quarenta anos.
\VS{37}Este é o Moisés que disse aos filhos de Israel: O Senhor, vosso Deus, vos levantará um profeta dentre vossos irmãos, como a mim; a ele ouvireis.
\VS{38}Este é aquele que esteve na congregação
{\ADD{do povo}} no deserto com o anjo que tinha lhe falado no monte Sinai, e
{\ADD{com}} nossos pais; o qual recebeu as palavras vivas, para dar a nós;
\VS{39}Ao qual nossos pais não quiseram obedecer; mas
{\ADD{o}} rejeitaram, e seus corações voltaram ao Egito;
\VS{40}Ao dizerem a Arão: Faz-nos deuses, que irão adiante de nós; porque
{\ADD{quanto a}} este Moisés, que nos levou para fora da terra do Egito, nós não sabemos o que aconteceu com ele.
\VS{41}E naqueles dias eles fizeram o bezerro, o ofereceram sacrifício ao ídolo, e se alegraram nas obras de suas
{\ADD{próprias}} mãos.
\VS{42}E Deus se afastou{\ADD{deles}} , e os entregou, para que servissem ao exército do céu, como está escrito no livro dos profetas: Ó casa de Israel, acaso foi a mim que oferecestes animais sacrificados, e ofertas no deserto por quarenta anos?
\VS{43}Porém tomastes
{\ADD{para si}} a tenda de Moloque, e a estrela de vosso Deus Renfã, figuras que vós fizeste para adorá-las; e eu
{\ADD{por isso}} vos expulsarei para além da Babilônia.
\VS{44}No deserto estava entre nossos pais o Tabernáculo do testemunho, assim como ele tinha ordenado, falando a Moisés, que o fizesse segundo o modelo que tinha visto.
\VS{45}O qual, recebendo
{\ADD{-o}} também nossos Pais, eles levaram com Josué para a possessão dos gentios que Deus expulsou diante de nossos Pais, até os dias de Davi;
\VS{46}O qual foi do agrado diante de Deus, e pediu para achar um tabernáculo para o Deus de Jacó.
\VS{47}E Salomão lhe construiu uma casa.
\VS{48}Mas o Altíssimo não habita em templos feitos por mãos, assim como o profeta diz:
\VS{49}O céu é o meu trono, e a terra
{\ADD{é}} o estrado dos meus pés; que casa vós construireis para mim?, diz o Senhor; ou Qual é o lugar do meu repouso?
\VS{50}Por acaso não
{\ADD{foi}} minhão mão
{\ADD{que}} fez todas estas coisas?
\VS{51}Vós, obstinados e incircuncisos de coração e de ouvidos! Vós sempre resistis ao Espírito Santo! Tal como vossos pais
{\ADD{foram}} , assim também
{\ADD{sois}} vós!
\VS{52}Qual dos profetas vossos pais não perseguiram? E eles mataram a todos os que anunciaram com antecedência a vinda do Justo, do qual agora vós tendes sido traidores e homicidas;
\VS{53}Que recebestes a Lei por ordem de anjos, e não
{\ADD{a}} guardastes.
\VS{54}Eles, ao ouvirem estas coisas, retalharam-se de raiva em seus corações, e rangiam os dentes contra ele.
\VS{55}Mas ele, estando cheio do Espírito Santo, olhando firmemente para o céu, viu à glória de Deus, e a Jesus, que estava à direita de Deus.
\VS{56}E disse: Eis que vejo os céus abertos, e o Filho do homem, que está a direita de Deus!
\VS{57}Mas eles, clamando com alta voz, taparam seus próprios ouvidos, e correram juntos contra ele;
\VS{58}E, lançando
{\ADD{-o}} fora da cidade,
{\ADD{o}} apedrejaram; e as testemunhas puseram as roupas deles junto aos pés de um rapaz chamado Saulo.
\VS{59}E apedrejaram a Estêvão, que estava clamando e dizendo: Senhor Jesus, recebe o meu espírito.
\VS{60}E pondo-se de joelhos, clamou com alta voz: Senhor, não os culpes por este pecado.E tendo dito isto, morreu.
\FTNT{*}{{\FR 7:60 }
morreu = lit. adormeceu
}

\par }\Chap{8}{\PP \VerseOne{1}E Saulo também consentia na morte dele. E naquele dia foi feita uma grande perseguição contra a igreja que estava em Jerusalém; e todos foram dispersos pelas regiões da Judeia, e de Samaria, exceto os apóstolos.
\VS{2}E
{\ADD{alguns}} homens devotos levaram juntos a Estêvão
{\ADD{para enterrá-lo}} ,e fizeram grande pranto por causa dele.
\VS{3}E Saulo tentava destruir a igreja, entrando nas casas, e puxando a homens e mulheres, entregava-os à prisão.
\VS{4}Os que, pois, estavam dispersos, passavam anunciando a palavra.
\VS{5}E Filipe, tendo descido à cidade de Samaria, pregava-lhes a Cristo.
\VS{6}E as multidões prestavam atenção em concordância às coisas que eram ditas por Filipe, ao ouvirem e verem os sinais que ele fazia.
\VS{7}Porque os espíritos imundos, clamando em alta voz, saíam de muitos que os tinham; e muitos paralíticos e aleijados foram curados.
\VS{8}E havia grande alegria naquela cidade.
\VS{9}E havia um certo homem, de nome Simão, que antes naquela cidade usava de magia, e fazia o povo de Samaria ficar admirado, dizendo de si mesmo ser alguém grande;
\VS{10}A quem eles todos davam atenção, desde o menor até o maior, diziam: Este é o grande poder de Deus.
\VS{11}E davam atenção a ele, porque com suas magias ele há muito tempo tinha lhes causado admiração.
\VS{12}Mas quando creram em Filipe, que lhes anunciava o Evangelho do Reino de Deus, e o nome de Jesus Cristo, eles foram batizados, tanto homens como mulheres.
\VS{13}E até mesmo Simão creu; e tendo sido batizado, ele continuou com Filipe; e vendo os sinais e maravilhas que eram feitas, ele ficou admirado.
\VS{14}E os apóstolos que estavam em Jerusalém, ao ouvirem que Samaria tinha recebido a palavra de Deus, enviaram-lhes a Pedro e a João.
\VS{15}Os quais, tendo descido, oraram por eles, para que recebessem ao Espírito Santo.
\VS{16}(Porque ainda sobre nenhum deles tinha descido; mas somente tinham sido batizados no nome do Senhor Jesus).
\VS{17}Então puseram as mãos sobre eles, e receberam o Espírito Santo.
\VS{18}E Simão ao ver que pela imposição das mãos dos apóstolos se dava o Espírito Santo, ofereceu-lhes dinheiro,
\VS{19}Dizendo: Dai também a mim este poder, que sobre qualquer um em quem eu puser as mãos, receba o Espírito Santo.
\VS{20}Mas Pedro lhe disse: Teu dinheiro seja contigo para perdição, porque pensaste que o dom de Deus
{\ADD{pudesse}} ser obtido por meio de dinheiro.
\VS{21}Tu não tens parte nem porção nesta palavra; porque teu coração não é correto diante de Deus.
\VS{22}Arrepende-te, pois, desta tua maldade, e ora a Deus, para que talvez este pensamento de teu coração te seja perdoado;
\VS{23}Porque eu vejo que tu estás em fel amargo, e atado em injustiça.
\VS{24}Mas respondendo Simão, disse: Orai vós por mim ao Senhor, para que nada do que dissestes venha sobre mim.
\VS{25}Tendo eles pois dado testemunho e falado a palavra do Senhor, voltaram a Jerusalém, e em muitas aldeias dos samaritanos anunciaram o Evangelho.
\VS{26}E um anjo do Senhor falou a Filipe, dizendo: Levanta-te, e vai para o sul, ao caminho que desce de Jerusalém para Gaza, que é deserto.
\VS{27}E ele levantou-se, e foi, e eis que um homem etíope, eunuco, administrador
{\ADD{subordinado}} a Candace, a rainha dos etíopes, que estava sobre
{\ADD{o controle}} de todos os bens dela, tinha vindo a Jerusalém para adorar;
\VS{28}E ele estava voltando, sentado em sua carruagem,
{\ADD{e}} lia ao profeta Isaías.
\VS{29}E o Espírito disse a Filipe: Aproxima-te, e ajunta-te a esta carruagem.
\VS{30}E Filipe, correndo, ouviu que ele estava lendo ao profeta Isaías, e disse: Tu entendes o que estás lendo?
\VS{31}E ele disse: Como eu poderia, se alguém não me ensinar? E pediu a Filipe que subisse e se sentasse com ele.
\VS{32}E o lugar da Escritura que ele estava lendo era este: Como ovelha ele foi levado ao matadouro, e como um cordeiro mudo fica diante do que o tosquia, assim
{\ADD{também}} ele não abriu sua boca.
\VS{33}Em sua humilhação foi tirado seu julgamento; e quem anunciará sua geração? Porque da terra sua vida é tirada.
\VS{34}E respondendo o eunuco a Filipe, disse: Eu te rogo, de quem o profeta diz isto? De si mesmo, ou de alguém outro?
\VS{35}E Filipe, abrindo sua boca, e começando desta escritura, anunciou-lhe o Evangelho
{\ADD{de}} Jesus.
\VS{36}E enquanto eles iam caminhando, chegaram a uma certa
{\ADD{porção}} de água; e o eunuco disse: Eis aqui água; o que me impede de ser batizado?
\VS{37}E Filipe disse: Se tu crês de todo coração,
{\ADD{então}} é lícito;E respondendo ele, disse: Creio que Jesus Cristo é o Filho de Deus.
\VS{38}E ele mandou parar a carruagem; e desceram ambos à água, tanto Filipe como o eunuco; e ele o batizou.
\VS{39}E quando eles subiram da água, o Espírito do Senhor arrebatou a Filipe, e o eunuco não mais o viu, porque ele foi
{\ADD{em}} seu caminho com alegria.
\VS{40}Mas Filipe se achou em Azoto; e passando, anunciava ao Evangelho
{\ADD{em}} todas as cidades, até que veio a Cesareia.

\par }\Chap{9}{\PP \VerseOne{1}E Saulo, ainda assoprando ameaças e mortes contra os discípulos do Senhor, foi ao chefe dos sacerdotes.
\VS{2}E pediu-lhe cartas para Damasco, para as sinagogas, para que se achasse alguns deste caminho, tanto homens como mulheres, ele
{\ADD{os}} trouxesse presos a Jerusalém.
\VS{3}E indo, aconteceu que chegando perto de Damasco, repentinamente brilhou ao redor dele uma luz do céu.
\VS{4}E caindo em terra, ouviu ma voz lhe dizendo: Saulo, Saulo, por que me persegues?
\VS{5}E ele disse: Quem és, Senhor?E o Senhor disse: Eu sou Jesus, a quem tu persegues; duro é para ti dar coices contra os aguilhões.
\VS{6}E ele, tremendo e atônito, disse: Senhor, que queres que eu faça?E o Senhor lhe
{\ADD{disse}} : Levanta-te, e entra na cidade; e
{\ADD{ali}} te será dito o que deves fazer.
\VS{7}E os homens que viajavam com ele pararam emudecidos, ouvindo, de fato, a voz, porém não vendo ninguém.
\VS{8}E Saulo se levantou da terra e, tendo aberto seus olhos, não via ninguém; e sendo guiado pela mão, levaram-no a Damasco.
\VS{9}E ele estava três dias sem ver; e não comeu, nem bebeu.
\VS{10}E havia em Damasco um certo discípulo, de nome Ananias; e disse-lhe o Senhor em visão: Ananias!E ele respondeu: Eis-me aqui, Senhor!
\VS{11}E o Senhor lhe
{\ADD{disse}} : Levanta-te, e vai a rua chamada Direita, e pergunta na casa de Judas por um chamado Saulo, de Tarso; porque que ele ora.
\VS{12}E ele viu em visão que um homem, de nome Ananias, entrava, e sobre ele punha a mão, para que voltasse a ver.
\VS{13}E Ananias respondeu: Senhor, eu ouvi de muitos sobre este homem, quantos males ele tem feito aos teus santos em Jerusalém;
\VS{14}E aqui ele tem poder dos chefes dos sacerdotes para prender a todos os que invocam o teu nome.
\VS{15}Mas o Senhor lhe disse: Vai, porque este me é um vaso escolhido para levar meu nome diante dos gentios, e dos reis, e dos filhos de Israel;
\VS{16}Porque eu mostrarei a ele o quanto ele deve sofrer por causa do meu nome.
\VS{17}E Ananias foi, e entrou na casa; e pondo as mãos sobre ele, disse: Irmão Saulo, o Senhor,
{\ADD{que é}} Jesus, aquele que apareceu a ti no caminho, me enviou para que tu voltes a ver, e sejas cheio do Espírito Santo.
\VS{18}E logo lhe caíram como que escamas dos olhos, e imediatamente voltou a ver; e levantando-se, foi batizado.
\VS{19}E ao comer, ele se fortaleceu. E Saulo ficou alguns dias com os discípulos
{\ADD{que estavam}} em Damasco.
\VS{20}E logo nas sinagogas pregava a Cristo,
{\ADD{dizendo}} que aquele era o Filho de Deus.
\VS{21}E todos os que o ouviam ficavam admirados, e diziam: Não é este aquele que em Jerusalém tentava destruir aos que invocavam este nome? E
{\ADD{não foi}} para isso
{\ADD{que}} ele veio aqui, para os levar presos aos chefes dos sacerdotes?
\VS{22}Mas Saulo se esforçava muito mais, e confundia aos judeus que habitavam em Damasco, provando que aquele era o Cristo.
\VS{23}E passados vários dias, os judeus tiveram conselho entre si para o matarem.
\VS{24}Mas as ciladas deles foram conhecidas por Saulo; e eles vigiavam as portas, tanto de dia como de noite, para poderem matá-lo.
\VS{25}Porém os discípulos, tomando-o de noite, levaram-no abaixo pelo muro em um cesto.
\VS{26}E Saulo, tendo vindo a Jerusalém, procurava se juntar aos discípulos; mas todos tinham medo dele, não crendo que fosse discípulo.
\VS{27}Mas Barnabás, tomando-o consigo, trouxe
{\ADD{-o}} aos apóstolos, e contou-lhes como no caminho tinha visto ao Senhor, e tinha lhe falado, e como em Damasco tinha falado ousadamente no nome de Jesus.
\VS{28}E ele estava junto deles, entrando e saindo em Jerusalém;
\VS{29}E falando ousadamente no nome do Senhor Jesus; falava e discutia também contra os gregos; mas eles procuravam matá-lo.
\VS{30}E os irmãos, ao perceberem
{\ADD{isto}} , o levaram até Cesareia, e o enviaram a Tarso.
\VS{31}Então as igrejas por toda a Judeia, e Galileia, e Samaria, tinham paz, e eram edificadas; e andando no temor do Senhor, e
{\ADD{na}} consolação do Espírito Santo, se multiplicavam.
\VS{32}E aconteceu que, Pedro, passando por todos
{\ADD{os lugares}} , veio também aos santos que habitavam em Lida.
\VS{33}E ali ele achou a um certo homem chamado Enéas, que havia oito anos que jazia numa cama, que era paralítico.
\VS{34}E Pedro lhe disse: Enéas, Jesus Cristo te cura; levanta-te, e faz tua cama. E logo ele se levantou.
\VS{35}E todos os que habitavam em Lida e Sarona o viram, os quais se converteram ao Senhor.
\VS{36}E havia em Jope uma certa discípula, de nome Tabita, que traduzido se diz Dorcas.
\FTNT{*}{{\FR 9:36 }
Dorcas = equiv. gazela
} Esta estava cheia de boas obras e doações que ela fazia
{\ADD{aos necessitados}} .
\VS{37}E aconteceu naqueles dias, que tendo ela ficado doente, morreu; e tendo a lavado, puseram-na no compartimento superior.
\VS{38}E como Lida era perto de Jope, os discípulos, ao ouvirem que Pedro estava ali, mandaram-lhe dois homens, rogando
{\ADD{-lhe}} que não demorasse a vir a eles.
\VS{39}E Pedro, tendo se levantado, foi com eles; o qual chegou, e
{\ADD{o}} levaram ao compartimento superior, e todas as viúvas o rodearam, chorando, e mostrando
{\ADD{-lhe}} as túnicas e roupas que Dorcas tinha feito quando estava com elas.
\VS{40}Mas Pedro, pondo para fora a todas; pôs-se de joelhos, e orou; e virando-se para o corpo, disse: Tabita, levanta-te;E ela abriu seus olhos, e vendo a Pedro, sentou-se.
\VS{41}E ele, dando-lhe a mão, levantou-a; e tendo chamado aos santos, e às viúvas, apresentou-a viva.
\VS{42}E isto ficou conhecido por toda Jope, e muitos creram no Senhor.
\VS{43}E aconteceu que ele ficou muitos dias em Jope, com um certo Simão curtidor.

\par }\Chap{10}{\PP \VerseOne{1}E havia um certo homem em Cesareia, de nome Cornélio, centurião, do esquadrão chamado Italiano;
\VS{2}Devoto, e temente a Deus, com toda a sua casa; e que fazia muitas doações ao povo, e continuamente orava a Deus.
\VS{3}Ele viu claramente em visão, cerca da hora nona do dia, a um anjo de Deus, que vinha a ele, e lhe dizia: Cornélio!
\VS{4}E ele, olhando-lhe atentamente, e muito atemorizado, disse: O que é, Senhor?E disse-lhe: Tuas orações e doações subiram à memória diante de Deus.
\VS{5}E agora envia homens a Jope, e manda chamar a Simão, que tem por sobrenome Pedro.
\VS{6}Este está hospedado na casa de um Simão curtidor, cuja casa é junto ao mar; este te dirá o que deves fazer.
\VS{7}E tendo partido o anjo que falava com Cornélio, ele chamou a dois de seus servos, e a um soldado devoto, dos que permaneciam continuamente com ele.
\VS{8}E tendo lhes contado tudo, enviou-os a Jope.
\VS{9}E no dia seguinte, enquanto estes iam pelo caminho, e chegando perto da cidade, Pedro subiu ao telhado para orar, quase à hora sexta.
\VS{10}E tendo ele fome, quis comer;
{\ADD{e}} enquanto estavam
{\ADD{lhe}} preparando, caiu sobre ele um êxtase.
\VS{11}E ele viu o céu aberto, e descia a ele um certo objeto, como um grande lençol, amarrado pelas quatro pontas, e abaixando-se à terra;
\VS{12}Em que havia de todos os animais quadrúpedes da terra, e animais selvagens, e répteis, e aves do céu.
\VS{13}E veio-lhe uma voz,
{\ADD{dizendo}} : Pedro, mata e come.
\VS{14}Mas Pedro disse: De maneira nenhuma, Senhor; porque nunca comi coisa alguma ordinária ou impura.
\VS{15}E a voz voltou a
{\ADD{dizer}} ,pela segunda vez: O que Deus purificou, não faças tu
{\ADD{como se fosse}} ordinário.
\VS{16}E isto aconteceu três vezes; e o objeto voltou a ser recolhido acima ao céu.
\VS{17}E enquanto Pedro estava pensando perplexo consigo mesmo o que seria aquela visão que ele tinha visto, eis que os homens que tinham sido enviados por Cornélio, perguntando pela casa de Simão, pararam à porta.
\VS{18}E chamando, perguntaram se Simão, que tinha por sobrenome Pedro, estava hospedado ali.
\VS{19}E estando Pedro pensando naquela visão, o Espírito lhe disse: Eis que três homens te buscam.
\VS{20}Então levanta-te, desce, e vai com eles, sem duvidar; porque eu os enviei.
\VS{21}E Pedro, tendo descido aos homens que tinham lhe sido enviados por Cornélio, disse: Eis que eu sou a quem buscais; qual é o motivo pelo qual estais aqui?
\VS{22}E eles disseram: Cornélio,
{\ADD{que é}} centurião, homem justo e temente a Deus,
{\ADD{e}} que tem
{\ADD{bom}} testemunho de toda a nação dos judeus, foi revelado por um santo anjo para te chamar até a casa dele, e ouvir de tuas palavras.
\VS{23}Então chamando-os para dentro, recebeu-os em casa. Mas no dia seguinte, Pedro foi com eles; e foram com ele alguns dos irmãos de Jope.
\VS{24}E no dia seguinte chegaram a Cesareia. E Cornélio estava esperando por eles, tendo chamado a seus parentes e amigos mais íntimos.
\VS{25}E sucedeu que, ao Pedro entrar, Cornélio se encontrou com ele, e caindo aos pés dele, adorou-o.
\VS{26}Mas Pedro o levantou, dizendo: Levanta-te; eu mesmo também sou um ser humano.
\VS{27}E tendo conversado com ele, entrou; e achou a muitos que
{\ADD{ali}} tinham se reunido.
\VS{28}E disse-lhes: Vós sabeis como não é lícito a um homem judeu juntar-se de estrangeiros, ou aproximar-se deles; mas Deus me mostrou que a ninguém chame de ordinário ou impuro.
\VS{29}Portanto eu, tendo sido chamado, vim sem qualquer oposição
{\ADD{de minha parte}} . Então eu pergunto: por que motivo me mandastes chamar?
\VS{30}E Cornélio disse: Há quatro dias que, até esta hora eu estava jejuando, e orava à hora nona em minha casa.
\VS{31}E eis que um homem se pôs diante de mim com uma roupa brilhante, e disse: Cornélio, tura oração tem sido ouvida, e tuas doações têm sido lembradas diante de Deus.
\VS{32}Envia pois a Jope, e manda chamar a Simão, que tem por sobrenome Pedro; este se hospeda na casa de Simão o curtidor, junto ao mar; quando ele vier, falará contigo.
\VS{33}Então logo eu enviei a ti; e bem fizeste em vir até aqui; agora pois estamos todos
{\ADD{aqui}} presentes diante de Deus, para ouvir tudo quanto Deus tem te ordenado.
\VS{34}E Pedro, abrindo a boca, disse: Reconheço que é verdade que Deus não faz acepção de pessoas.
\VS{35}Mas sim, em toda nação, aquele que o teme, e pratica a justiça,
{\ADD{este}} lhe é agradável.
\VS{36}A palavra que Deus enviou aos filhos de Israel, anunciando o Evangelho da paz por meio de Jesus Cristo; este é o Senhor de todos.
\VS{37}Vós sabeis da palavra que veio por toda a Judeia, começando desde a Galileia, depois do batismo que João pregou.
\VS{38}{\ADD{E}} sobre Jesus de Nazaré; como Deus o ungiu com o Espírito Santo, e com poder; o qual percorreu
{\ADD{os lugares}} fazendo o bem, e curando a todos os oprimidos pelo diabo; porque Deus era com ele.
\VS{39}E nós somos testemunhas de todas as coisas que ele fez; tanto na terra dos judeus, como em Jerusalém; ao qual mataram, pendurando
{\ADD{-o}} em um madeiro.
\VS{40}A este Deus ressuscitou ao terceiro dia, e fez com que fosse manifesto;
\VS{41}Não a todo o povo, mas sim a testemunhas determinadas por Deus com antecedência: a nós, que juntamente com ele comemos e bebemos, depois dele ter ressuscitado dos mortos.
\VS{42}E ele nos mandou pregar ao povo, e dar testemunho de que ele é o que foi ordenado por Deus
{\ADD{para ser}} Juiz dos vivos e dos mortos.
\VS{43}A este todos os profetas dão testemunho, de que todos os que nele crerem receberão perdão dos pecados por meio do seu nome.
\VS{44}E estando Pedro ainda falando estas palavras, o Espírito Santo caiu sobre todos os que ouviam a palavra.
\VS{45}E os crentes que eram da circuncisão, tantos quantos tinham vindo com Pedro, ficaram muito admirados de que também sobre os gentios fosse derramado o dom do Espírito Santo.
\VS{46}Porque eles os ouviam falar em
{\ADD{diversas}} línguas, e a engrandecer a Deus. Então Pedro respondeu:
\VS{47}Por acaso pode alguém impedir a água, para que não sejam batizados estes, que também, assim como nós, receberam o Espírito Santo?
\VS{48}E mandou que fossem batizados no nome do Senhor. Então lhe pediram que continuasse
{\ADD{com eles}} por alguns dias.

\par }\Chap{11}{\PP \VerseOne{1}E os apóstolos, e os irmãos que estavam na Judeia, ouviram que também os gentios receberam a palavra de Deus.
\VS{2}E quando Pedro subiu a Jerusalém, discutiam contra ele os que eram da circuncisão;
\VS{3}Dizendo: Tu entraste
{\ADD{na casa}} de homens incircuncisos, e comeste com eles.
\VS{4}Mas Pedro começou a lhes explicar
{\ADD{tudo}} em ordem, dizendo:
\VS{5}Eu estava orando na cidade de Jope, e vi em êxtase uma visão: um certo objeto que descia como um grande lençol, pelas quatro pontas abaixado desde o céu, e vinha até mim.
\VS{6}No qual, olhando eu com atenção, considerei e vi quatro quadrúpedes da terra, e animais selvagens, e répteis, e aves do céu.
\VS{7}E eu ouvi uma voz que me dizia: Levanta-te Pedro, mata e come.
\VS{8}Mas eu disse: De maneira nenhuma, Senhor; porque nunca comi coisa alguma ordinária, nem coisa imunda entrou em minha boca.
\VS{9}E a voz me respondeu do céu pela segunda vez: O que Deus purificou, não o faças tu
{\ADD{como}} ordinário.
\VS{10}E isto aconteceu por três vezes; e voltou-se tudo a recolher acima ao céu.
\VS{11}E eis que logo três homens, enviados a mim de Cesareia, pararam junto à casa onde eu estava.
\VS{12}E o Espírito me disse que eu fosse com eles, sem duvidar; e também estes seis irmãos foram comigo, e entramos na casa daquele homem.
\VS{13}E ele nos contou como tinha visto um anjo estar em sua casa, e tinha lhe dito: Envia
{\ADD{alguns}} homens a Jope, e manda chamar a Simão, que tem por sobrenome Pedro;
\VS{14}O qual te falará palavras, em que tu sejas salvo, e
{\ADD{também}} toda a tua casa.
\VS{15}E quando comecei a falar, caiu o Espírito Santo sobre eles, assim como também no princípio
{\ADD{tinha caído}} sobre nós.
\VS{16}E eu me lembrei da palavra do Senhor, como ele tinha dito: Verdadeiramente João batizou com água, mas vós sereis batizados com o Espírito Santo.
\VS{17}Portanto, se Deus deu a eles igual dom, assim como também a nós, que temos crido no Senhor Jesus Cristo; quem era eu, pois, para que pudesse proibir a Deus?
\VS{18}E ao ouvirem estas coisas, se acalmaram, e glorificavam a Deus, dizendo: Portanto também aos gentios Deus deu arrependimento para a vida.
\VS{19}E os que foram dispersos por causa da perseguição que aconteceu por causa de Estêvão, passaram até a Fenícia, e Chipre, e Antioquia; não falando a ninguém a palavra, a não ser somente aos judeus.
\VS{20}E havia deles alguns homens cipriotas e cirenenses, os quais ao entrarem em Antioquia, falaram aos gregos, anunciando o Evangelho do Senhor Jesus.
\VS{21}E a mão do Senhor era com eles, e um grande número creu, e se converteu ao Senhor.
\VS{22}E esta notícia sobre eles chegou aos ouvidos da igreja que estava em Jerusalém; e enviaram a Barnabé, para ir até Antioquia.
\VS{23}O qual, ao chegar
{\ADD{lá}} ,e tendo visto a graça de Deus, alegrou-se; e exortou a todos, para que com o propósito do coração permanecessem no Senhor.
\VS{24}Porque ele era um bom homem, e cheio do Espírito Santo, e de fé; e uma grande multidão foi acrescentada ao Senhor.
\VS{25}E Barnabé foi para Tarso, para buscar a Saulo; e achando-o, trouxe-o a Antioquia.
\VS{26}E sucedeu que,
{\ADD{durante}} um ano completo eles se congregaram naquela igreja, e ensinavam a uma grande multidão; e em Antioquia os discípulos foram chamados pela primeira vez de cristãos.
\VS{27}E naqueles dias desceram de Jerusalém
{\ADD{alguns}} profetas a Antioquia.
\VS{28}E levantando-se um deles, por nome Ágabo, declarou pelo Espírito, que estava para haver uma grande fome em todo o mundo; que veio a acontecer no tempo de Cláudio César.
\VS{29}E os discípulos determinaram de cada um, conforme o que pudesse, mandar algum
{\ADD{socorro}} para serviço dos irmãos que habitavam na Judeia.
\VS{30}O que também fizeram, enviando
{\ADD{-o}} aos anciãos
\FTNT{*}{{\FR 11:30 }
anciãos = equiv. presbíteros
} pela mão de Barnabé e de Saulo.

\par }\Chap{12}{\PP \VerseOne{1}E por aquele mesmo tempo o rei Herodes pôs as mãos para maltratar a alguns da igreja.
\VS{2}E matou a Tiago, o irmão de João, pela espada.
\VS{3}E vendo que isto agradava aos judeus, ele fez ainda mais, para também prender a Pedro (e eram os dias dos
{\ADD{pães}} sem fermento).
\VS{4}Do qual, também detendo, lançou-o na prisão, entregando
{\ADD{-o}} a quatro quaternos de soldados, que o guardassem; pretendendo tirá-lo
{\ADD{para mostrá-lo}} ao povo depois da Páscoa.
\VS{5}Então Pedro era mantido na prisão; mas a igreja fazia fervorosa oração a Deus por ele.
\VS{6}E quando Herodes estava para tirá-lo
{\ADD{para apresentá-lo}} ,naquela mesma noite Pedro estava dormindo entre dois soldados, acorrentado com duas correntes; e os guardas diante da porta guardavam a prisão.
\VS{7}E eis que veio acima um anjo do Senhor, e uma luz brilhou na prisão; e tocando em Pedro em sua lateral, despertou-o, dizendo; Levanta-te, depressa!E as correntes caíram de suas mãos.
\VS{8}E o anjo lhe disse: Arruma-te, e amarra as tuas sandálias.E ele fez assim. E disse-lhe: Põe tua capa sobre ti, e segue-me.
\VS{9}E saindo, o seguia; e não sabia que era verdade o que se fazia pelo anjo, mas pensava que estava tendo alguma visão.
\VS{10}E ao passarem a primeira e a segunda guarda, chegaram à porta de ferro, que leva à cidade, a qual foi aberta por si mesma; e tendo saído, foram a uma rua, e logo o anjo partiu dele.
\VS{11}E tendo Pedro voltado a si, disse: Agora eu sei verdadeiramente que o Senhor enviou a seu anjo, e me livrou da mão de Herodes, e de toda a expectativa do povo dos judeus.
\VS{12}E ele, reconhecendo
{\ADD{isto}} ,foi à casa de Maria, a mãe de João, que tinha por sobrenome Marcos, onde muitos estavam juntos, e oravam.
\VS{13}E Pedro, tendo batido a porta da entrada,veio uma moça de nome Rode, para escutar.
\VS{14}E ela, reconhecendo a voz de Pedro, de alegria não abriu a porta da entrada, em vez disso ela correu para dentro, e anunciou que Pedro estava fora à porta da entrada.
\VS{15}E lhe disseram: Tu estás delirando.Mas ela, insistindo que assim era. E eles diziam: É o anjo dele.
\VS{16}Mas Pedro continuava a bater; e ao abrirem, viram-no, e ficaram espantados.
\VS{17}Mas ele, fazendo-lhes gestos para que calassem, contou-lhes como o Senhor tinha lhe tirado da prisão, e disse: Anunciai isto a Tiago e aos irmãos.E tendo saído, foi para outro lugar.
\VS{18}E vindo o dia, havia não pouca perturbação entre os soldados, sobre o que, pois, tinha acontecido com Pedro.
\VS{19}E quando Herodes o buscou, e não o achou, tendo investigado aos guardas, mandou que eles fossem levados
{\ADD{para serem mortos}} . E partindo da Judeia para Cesareia, ficou
{\ADD{ali}} .
\VS{20}E Herodes estava extremamente irritado com os de Tiro e de Sidom; porém eles, vindo em concordância até ele, e persuadindo a Blasto, que era o camareiro do rei, pediram paz, porque a terra deles dependia dos alimentos da terra do rei
{\ADD{Herodes}} .
\VS{21}E num dia marcado, Herodes vestiu roupas reais, e sentando no tribunal, fez-lhes um discurso.
\VS{22}E o povo exclamava: Voz de deus, e não de homem!
\VS{23}E no mesmo instante um anjo do Senhor o feriu, porque ele não deu a glória a Deus; e tendo sido comido por vermes, deixou de respirar.
\VS{24}E a palavra de Deus crescia, e se multiplicava.
\VS{25}E Barnabé e Saulo, tendo cumprido aquele serviço, voltaram a Jerusalém, tomando também consigo a João, o que tinha por sobrenome Marcos.

\par }\Chap{13}{\PP \VerseOne{1}E havia em Antioquia, na igreja que estava
{\ADD{ali}} , alguns profetas e mestres: Barnabé e Simeão, chamado Níger, e Lúcio cireneu, e Manaem, que tinha sido criado na infância junto com Herodes o Tetrarca, e Saulo.
\VS{2}E tendo eles prestado serviço ao Senhor, e jejuado, o Espirito Santo disse: Separai-me a Barnabé e a Saulo, para a obra para a qual eu os tenho chamado.
\VS{3}Então jejuando, e orando, e pondo as mãos sobre eles,
{\ADD{os}} despediram.
\VS{4}Portanto estes, enviados pelo Espírito Santo, desceram a Selêucia, e dali navegaram para o Chipre.
\VS{5}E tendo chegado a Salamina, anunciavam a palavra de Deus nas sinagogas dos judeus; e também tinham a João como trabalhador
{\ADD{para os auxiliar}} .
\VS{6}E tendo eles atravessado a Ilha até Pafo, acharam a um certo mago, falso profeta, judeu, cujo nome era Barjesus.
\VS{7}O qual estava com o procônsul Sérgio Paulo, homem prudente. Este, tendo chamado a si a Barnabé, e a Saulo, procurava muito ouvir a palavra de Deus.
\VS{8}Mas resistia-lhes Elimas, o mago (que assim significa seu nome), procurando afastar o procônsul da fé.
\VS{9}Mas Saulo, que também
{\ADD{se chama}} Paulo, cheio do Espírito Santo, e olhando fixamente para ele, disse:
\VS{10}Ó filho do diabo, cheio de toda enganação e toda malícia, inimigo de toda justiça, não cessarás de perverter os corretos caminhos do Senhor?
\VS{11}E agora, eis que a mão do Senhor
{\ADD{está}} contra ti, e serás cego, não vendo o sol por algum tempo.E no mesmo instante caiu sobre ele um embaçamento, e trevas; e ele, andando ao redor, procurava alguém que o guiasse pela mão.
\VS{12}Então o procônsul, vendo o que tinha acontecido, creu, espantado pela doutrina do Senhor.
\VS{13}E tendo partido de Pafo, Paulo e os que estavam com ele foram a Perges,
{\ADD{cidade}} da Panfília. Mas João, separando-se deles, voltou a Jerusalém.
\VS{14}E eles, tendo passado de Perges, vieram a Antioquia,
{\ADD{cidade}} da Pisídia; e ao entrarem na sinagoga
{\ADD{n}} um dia de sábado, sentaram-se.
\VS{15}E depois da leitura da Lei e dos profetas, os chefes da sinagoga lhes mandaram, dizendo: Homens irmãos, se em vós há
{\ADD{alguma}} palavra de exortação ao povo, dizei.
\VS{16}E Paulo, levantando-se e fazendo gesto com a mão, disse: Homens israelitas, e
{\ADD{vós}} os que temeis a Deus, ouvi:
\VS{17}O Deus deste povo de Israel escolheu a nossos pais, e exaltou ao povo, sendo eles estrangeiros na terra do Egito, e com o braço levantando, ele os tirou dela.
\VS{18}E pelo tempo de cerca de quarenta anos, ele suportou os costumes deles no deserto.
\VS{19}E tendo destruído a sete nações na terra de Canaã, repartiu-lhes as terras por sorte.
\VS{20}E depois disto, cerca de quatrocentos e cinquenta anos, ele
{\ADD{lhes}} deu juízes, até o profeta Samuel.
\VS{21}E depois disto, pediram a um Rei, e ele lhes deu a Saul, filho de Quis, homem da tribo de Benjamim,
{\ADD{durante}} quarenta anos.
\VS{22}E tirando a este, levantou-lhes por rei a Davi, ao qual também deu testemunho, e disse: Eu achei a Davi,
{\ADD{filho}} de Jessé, um homem conforme o meu coração, que fará toda a minha vontade.
\VS{23}Da descendência
\FTNT{*}{{\FR 13:23 }
lit. semente
} deste, conforme a promessa, Deus levantou a Jesus por Salvador de Israel;
\VS{24}Tendo João primeiro, antes de sua vinda, pregado o batismo de arrependimento a todo o povo de Israel.
\VS{25}Mas quando João cumpriu
{\ADD{sua}} carreira, disse: Quem vós pensais que eu sou? Eu não sou o
{\ADD{Cristo}} ,mas eis que após mim vem aquele, cujas sandálias dos pés eu não sou digno de desatar.
\VS{26}Homens irmãos, filhos da descendência de Abraão, e os que dentre vós temem a Deus, a vós é enviada a palavra desta salvação.
\VS{27}Porque os que habitavam em Jerusalém, e seus líderes, não conhecendo a este, ao condenarem
{\ADD{-no}} ,cumpriram as vozes dos profetas, que são lidas todos os sábados.
\VS{28}E
{\ADD{mesmo}} achando nenhum motivo para morte, pediram a Pilatos que fosse morto.
\VS{29}E tendo eles cumprido todas as coisas que estavam escritas sobre ele, tirando
{\ADD{-o}} do madeiro, puseram
{\ADD{-no}} na sepultura.
\VS{30}Mas Deus o ressuscitou dos mortos.
\VS{31}O qual foi visto durante muitos dias pelos que haviam subido com ele da Galileia, que são suas testemunhas para com o povo.
\VS{32}E nós vos anunciamos o Evangelho da promessa que foi feita aos pais; ao qual Deus já nos cumpriu a nós, filhos deles, ressuscitando a Jesus.
\VS{33}Assim como está escrito no salmo segundo: Tu és meu Filho, hoje eu te gerei.
\VS{34}E
{\ADD{quanto a}} que o ressuscitasse dos mortos, para nunca mais voltar à degradação,
\FTNT{†}{{\FR 13:34 }
Ou: deterioração, putrefação. Também nos versículos seguintes.
} assim disse: Eu vos darei as fiéis beneficências de Davi.
\VS{35}Por isso que também em outro
{\ADD{salmo}} ele diz: Não permitirás que teu Santo veja degradação.
\VS{36}Porque, na verdade, tendo Davi servido ao conselho de Deus, morreu,
\FTNT{‡}{{\FR 13:36 }
morreu = lit. adormeceu
} foi posto junto a seus pais, e viu degradação.
\VS{37}Mas aquele a quem Deus ressuscitou nenhuma degradação viu.
\VS{38}Seja-vos pois conhecido, homens irmãos, que por este vos é anunciado o perdão dos pecados.
\VS{39}E de tudo o que pela lei de Moisés não pudestes ser justificados, neste é justificado todo aquele que crê.
\VS{40}Então vede, para que não venha sobre vós o que está escrito nos
{\ADD{livros dos}} profetas:
\VS{41}{\ADD{Vós}} desprezadores, vede e espantai-vos, e desaparecei-vos; porque eu opero obra em vossos dias, obra na qual não crereis, se alguém vos contar.
\VS{42}E tendo os judeus saído da sinagoga,os gentios rogaram
{\ADD{-lhes}} que no sábado seguinte eles lhes falassem estas palavras.
\VS{43}E tendo terminado
{\ADD{a reunião}} da sinagoga, muitos dos judeus, e dos religiosos prosélitos, seguiram a Paulo e a Barnabé; os quais, falando-lhes, exortavam-nos a permanecerem na graça de Deus.
\VS{44}E no sábado seguinte ajuntou-se quase toda a cidade para ouvir a palavra de Deus.
\VS{45}Mas os Judeus, ao verem as multidões, ficaram cheios de inveja, e falavam contrariamente ao que Paulo dizia, falando contrariamente e blasfemando.
\VS{46}Mas Paulo e Barnabé, usando de ousadia, disseram: Era necessário que a palavra de Deus fosse primeiro falada a vós; mas já que a rejeitais, e não vos julgais dignos da vida eterna, eis que nós nos viramos
{\ADD{em direção}} aos gentios.
\VS{47}Porque assim o Senhor nos mandou,
{\ADD{dizendo}} : Eu te pus como luz para os gentios, para que tu sejas como salvação até às extremidades da terra.
\VS{48}E os gentios, tendo ouvido
{\ADD{isto}} ,alegraram-se, e glorificavam ao Senhor; e creram todos quantos estavam determinados para a vida eterna.
\VS{49}E a palavra do Senhor era divulgada por toda aquela região.
\VS{50}Mas os judeus incitaram algumas mulheres devotas e honradas, e aos líderes da cidade, e levantaram perseguição contra Paulo e Barnabé, e os expulsaram de seus limites.
\VS{51}Mas eles, sacudindo contra eles o pó dos seus pés, vieram a Icônio.
\VS{52}E os discípulos se enchiam de alegria e do Espírito Santo.

\par }\Chap{14}{\PP \VerseOne{1}E aconteceu em Icônio, que entraram juntos na sinagoga dos judeus, e falaram de tal maneira que creu uma grande multidão, tanto de judeus como de gregos.
\VS{2}Mas os judeus incrédulos incitaram e irritaram os ânimos dos gentios contra os irmãos.
\VS{3}Então eles ficaram
{\ADD{ali}} por muito tempo, falando ousadamente no Senhor, o qual dava testemunho à palavra de sua graça, concedendo
{\ADD{que}} sinais e milagres fossem feitos pelas mãos deles.
\VS{4}E a multidão da cidade se dividiu; e uns eram a favor dos judeus, e outros a favor dos apóstolos.
\VS{5}E fazendo-se uma rebelião, tanto de judeus como de gentios, juntos com seus líderes, para falarem mal deles, e os apedrejarem.
\VS{6}E eles, sabendo
{\ADD{disto}} ,fugiram para as cidades da Licaônia,
{\ADD{chamadas}} Listra e Derbe; e à região ao redor.
\VS{7}E ali eles anunciavam ao Evangelho.
\VS{8}E um certo homem em Listra estava sentado, tendo incapacidade nos pés, aleijado desde o ventre de sua mãe, que nunca tinha andado.
\VS{9}Este ouviu Paulo falando; o qual, olhando com atenção, e vendo que ele tinha fé para ser curado,
\VS{10}Disse em alta voz: Levanta-te direito sobre teus pés.E ele saltou e andou.
\VS{11}E as multidões, vendo o que Paulo tinha feito, levantaram suas vozes, dizendo em
{\ADD{língua}} licaônica: Os deuses se fizeram semelhantes a homens, e desceram até nós.
\VS{12}E chamaram a Barnabé de Júpiter; e a Paulo, de Mercúrio; porque este era o líder ao falar.
\VS{13}E o sacerdote de Júpiter, que estava diante da cidade deles, trazendo touros e grinaldas à entrada da porta, ele, junto com as companhias, queria oferecer sacrifício
{\ADD{a eles}} .
\VS{14}Mas os apóstolos Barnabé e Paulo, ao ouvirem
{\ADD{isto}} ,rasgaram suas roupas, e saltaram entre a multidão, clamando,
\VS{15}E dizendo: Homens, por que fazeis estais coisas? Também nós somos homens como vós, sujeitos às mesmas emoções; e vos anunciamos o Evangelho para que vos convertais destas vaidades para o Deus vivo, que fez o céu, a terra, o mar, e tudo quanto neles há.
\VS{16}O qual nas gerações passadas deixou os os gentios andarem seus
{\ADD{próprios}} caminhos.
\VS{17}Ainda que, contudo, não tenha deixado a si mesmo sem testemunho, fazendo o bem desde o céu, dando-nos chuvas, e tempos frutíferos,
{\ADD{e}} enchendo nossos corações de alimento e alegria.
\VS{18}E tendo disto isto, apenas detiveram as multidões de que não fizessem sacrifícios a eles.
\VS{19}Mas vieram
{\ADD{alguns}} judeus de Antioquia, e de Icônio, e persuadiram a multidão; e apedrejando a Paulo, arrastaram
{\ADD{-no}} para fora da cidade, pensando que ele estivesse morto.
\VS{20}Mas, tendo os discípulos ficado ao seu redor, ele se levantou, e entrou na cidade; e no dia seguinte saiu com Barnabé para Derbe.
\VS{21}E tendo anunciado o Evangelho àquela cidade, e feito muitos discípulos, eles voltaram a Listra, e a Icônio, e a Antioquia,
\VS{22}Confirmando os ânimos dos discípulos, e exortando-os para que permanecessem na fé, e que nos é necessário entrar no Reino de Deus por meio de muitas aflições.
\VS{23}E tendo escolhido por votação anciãos
\FTNT{*}{{\FR 14:23 }
anciãos = equiv. presbíteros
} para cada igreja, orando com jejuns, eles foram enviados ao Senhor, no qual tinham crido.
\VS{24}E tendo passado por Pisídia, vieram à Panfília.
\VS{25}E tendo falado a palavra em Perges, desceram a Atália.
\VS{26}E dali navegaram para Antioquia, onde tinham sido encomendados à graça de Deus para a obra que eles
{\ADD{já}} tinham cumprido.
\VS{27}E ao chegarem, e reunirem a igreja, relataram quão grandes coisas Deus tinha feito com eles; e como ele tinha aberto a porta da fé aos gentios.
\VS{28}E eles ficaram ali não pouco tempo com os discípulos.

\par }\Chap{15}{\PP \VerseOne{1}E alguns que tinham descido da Judeia ensinavam aos irmãos,
{\ADD{dizendo}} : Se vós não vos circuncidardes conforme o costume de Moisés, não podeis ser salvos.
\VS{2}Então, havendo não pequena resistência e confronto de Paulo e Barnabé contra eles, ordenaram que Paulo, Barnabé e alguns outros deles subissem aos apóstolos e aos anciãos a Jerusalém sobre esta questão.
\VS{3}Então sendo eles preparados para a viagem e despedidos pela igreja, passaram pela Fenícia e Samaria, contando
{\ADD{sobre}} a conversão dos gentios; e davam grande alegria a todos os irmãos.
\VS{4}E tendo chegado a Jerusalém, eles foram recebidos pela igreja, e pelos apóstolos e anciãos; e lhes anunciaram quão grandes coisas Deus tinha feito com eles;
\VS{5}Mas
{\ADD{que}} alguns do grupo dos fariseus que tinham crido, levantaram-se, dizendo que era necessário circuncidá-los, e mandar
{\ADD{-lhes}} que guardem a Lei de Moisés.
\VS{6}E os apóstolos e anciãos se reuniram para dar atenção a este assunto.
\VS{7}E havendo muita discussão, Pedro se levantou, e lhes disse: Homens irmãos, vós sabeis que há muito tempo Deus
{\ADD{meu}} escolheu entre nós, para que por minha boca os gentios ouvissem a palavra do Evangelho, e cressem.
\VS{8}E Deus, que conhece os corações, deu-lhes testemunho, dando-lhes o Espírito Santo, assim como também a nós.
\VS{9}E nenhuma diferença fez entre nós e eles, purificando seus corações pela fé.
\VS{10}Então agora, por que tentais a Deus, pondo um jugo sobre o pescoço dos discípulos; que nem nossos pais, nem nós podemos levar?
\VS{11}Mas cremos que, pela graça do Senhor Jesus Cristo, nós somos salvos, assim como também eles.
\VS{12}E toda a multidão se calou; e ouviram a Barnabé e a Paulo, que contavam quão grandes sinais e milagres Deus tinha feito por meio deles entre os gentios.
\VS{13}E tendo estes se calado, Tiago respondeu, dizendo: Homens irmãos, ouvi-me:
\VS{14}Simão informou como primeiro Deus visitou aos gentios, para tomar
{\ADD{deles}} um povo para seu nome.
\VS{15}E com isso concordam as palavras dos profetas, como está escrito:
\VS{16}Depois disto eu voltarei, e reconstruirei o tabernáculo de Davi, que caído está; e reconstruirei
{\ADD{de}} suas ruínas, e voltarei a levantá-lo;
\VS{17}Para que o resto da humanidade busque ao Senhor, e todos os gentios
\FTNT{*}{{\FR 15:17 }
gentios
equiv. nações
} sobre os quais meu nome é invocado, diz o Senhor, que faz todas estas coisas.
\VS{18}São conhecidas por Deus desde a antiguidade todas as suas obras.
\VS{19}Portanto eu julgo que aqueles que dos gentios se convertem a Deus não devem ser perturbados.
\VS{20}Mas
{\ADD{que}} lhes escrevamos para que se abstenham das contaminações dos ídolos, e do pecado sexual, e da
{\ADD{carne}} sufocada, e do sangue.
\VS{21}Porque Moisés, desde as gerações antigas, tem em cada cidade quem o pregue nas sinagogas, sendo lido todo sábado.
\VS{22}Então pareceu bem aos apóstolos, e aos anciãos, com toda a igreja, eleger deles
{\ADD{alguns}} homens, para serem enviados com Paulo e Barnabé a Antioquia: Judas, que tinha por sobrenome Barsabás; e a Silas, homens líderes entre os irmãos.
\VS{23}E escreveram por meio deles o seguinte: Os apóstolos e os anciãos, e os irmãos – para os irmãos dentre os gentios, que
{\ADD{estão}} em Antioquia, Síria e Cilícia; saudações.
\VS{24}Dado que ouvimos que alguns dos que saíram de nós vos perturbaram com palavras,
{\ADD{e}} causaram incômodo a vossas almas, dizendo que deveis vos circuncidar e guardar a Lei, aos quais não mandamos;
\VS{25}Pareceu-nos bem, reunidos em concordância, escolher
{\ADD{alguns}} homens, e enviá-los até vós, com nossos amados Barnabé e Paulo.
\VS{26}Homens que têm arriscado suas vidas pelo nome de nosso Senhor Jesus Cristo.
\VS{27}Então enviamos a Judas e a Silas, os quais
{\ADD{vos}} dirão as mesmas coisas pessoalmente.
\VS{28}Porque pareceu bem ao Espírito Santo e a nós, de nenhuma carga a mais vos impor, a não ser estas coisas necessárias:
\VS{29}Que vos abstenhais das coisas sacrificadas aos ídolos, do sangue, da
{\ADD{carne}} sufocada, e do pecado sexual; das quais, se vos guardardes, fareis bem. Que o bem vos suceda.
\VS{30}Sendo, pois, eles despedidos, vieram a Antioquia, e reunindo a multidão, entregaram a carta.
\VS{31}E ao lerem, alegraram-se pela consolação.
\VS{32}E então Judas e Silas, sendo também profetas, com muitas palavras exortaram e firmaram aos irmãos.
\VS{33}E ficando
{\ADD{ali}} por algum tempo, permitiram que voltassem em paz dos irmãos para os apóstolos.
\VS{34}Mas a Silas pareceu bem continuar ali.
\VS{35}E Paulo e Barnabé ficaram em Antioquia, ensinando e evangelizando, com também muitos outros, a palavra do Senhor.
\VS{36}E depois de alguns dias, Paulo disse a Barnabé: Voltemos a visitar a nossos irmãos em cada cidade onde tenhamos anunciado a palavra do Senhor,
{\ADD{para ver}} como estão.
\VS{37}E Barnabé aconselhou para que tomassem consigo a João, chamado Marcos.
\VS{38}Mas Paulo achou adequado que não tomassem consigo a aquele que desde a Panfília tinha se separado deles, e não tinha ido com eles para
{\ADD{aquela}} obra.
\VS{39}Houve então
{\ADD{entre eles}} tal discórdia, que eles se separaram um do outro; e Barnabé, tomando consigo a Marcos, navegou para o Chipre.
\VS{40}Mas Paulo, escolhendo a Silas, partiu-se, enviado pelos irmãos para a graça de Deus.
\VS{41}E ele passou pela Síria e Cilícia, firmando as igrejas.

\par }\Chap{16}{\PP \VerseOne{1}E ele veio a Derbe e Listra; e eis que estava ali um certo discípulo, de nome Timóteo, filho de uma certa mulher judia crente, mas de pai grego.
\VS{2}Do qual era
{\ADD{bem}} testemunhado pelos irmãos
{\ADD{que estavam}} em Listra e Icônio.
\VS{3}A este Paulo quis que fosse com ele; e tomando-o, circuncidou-o, por causa dos judeus, que estavam naqueles lugares; porque todos conheciam o pai dele, que era grego.
\VS{4}E eles, passando pelas cidades, entregavam-lhes as ordenanças que foram determinadas pelos apóstolos e anciãos em Jerusalém, para que
{\ADD{as}} guardassem.
\VS{5}E assim as Igrejas eram firmadas na fé, e cada dia aumentavam em número.
\VS{6}E passando pela Frígia, e pela região da Galácia, foi-lhes impedido pelo Espírito Santo de falarem a palavra na Ásia.
\VS{7}{\ADD{E}} quando eles vieram a Mísia, tentaram ir à Bitínia; mas o Espírito não lhes permitiu.
\VS{8}E tendo passado por Mísia, desceram a Trôade.
\VS{9}E uma visão foi vista por Paulo durante a noite: um homem Macedônio se pôs
{\ADD{diante dele}} ,rogando-lhe, e dizendo: Passa à Macedônia, e ajuda-nos!
\VS{10}E quando ele viu a visão, logo procuramos partir para a Macedônia, concluindo que o Senhor estava nos chamando para anunciarmos o Evangelho a eles.
\VS{11}Então, tendo navegado desde Trôade, viemos correndo caminho direto a Samotrácia, e no
{\ADD{dia}} seguinte a Neápolis.
\VS{12}E dali a Filipos, que é a primeira cidade desta parte da Macedônia,
{\ADD{e é}} uma colônia; e estivemos naquela cidade
{\ADD{por}} alguns dias.
\VS{13}E no dia de sábado saímos para fora da cidade, onde costumava ser feita oração; e tendo
{\ADD{nos}} sentado, falamos às mulheres que tinham se ajuntado
{\ADD{ali}} .
\VS{14}E uma certa mulher, por nome Lídia, vendedora de púrpura, da cidade de Tiatira, que servia a Deus;
{\ADD{ela nos}} ouviu; o coração da qual o Senhor abriu, para que prestasse atenção ao que Paulo dizia.
\VS{15}E quando ela foi batizada, e
{\ADD{também}} sua casa, ela
{\ADD{nos}} rogou, dizendo: Se vós tendes julgado que eu seja fiel ao Senhor, entrai em minha casa e ficai ali.E ela insistiu para conosco.
\VS{16}E aconteceu, que ao estarmos nós indo à oração, saiu ao nosso encontro uma moça que tinha espírito de pitonisa; a qual ao fazer adivinhações trazia grande lucro a seus senhores.
\VS{17}Esta, seguindo após Paulo e nós, clamava, dizendo: Estes homens são servos do Deus Altíssimo, que nos anunciam o caminho da salvação.
\VS{18}E ela fazia isto por muitos dias. Mas Paulo, estando descontente com isto, virou-se, e disse ao espírito: Em nome de Jesus Cristo eu te mando que saias dela.E na mesma hora
{\ADD{o espírito}} saiu.
\VS{19}E os senhores dela, vendo que a esperança de lucro deles tinha ido embora, pegaram a Paulo e a Silas, e
{\ADD{os}} levaram à praça, diante dos governantes.
\VS{20}E apresentando-os aos magistrados, disseram: Estes homens perturbam nossa cidade, sendo judeus;
\VS{21}E eles anunciam costumes que não nos é lícito receber, nem fazer; pois somos romanos.
\VS{22}E a multidão se levantou juntamente contra eles; e os oficiais, rasgando suas roupas, mandaram que fossem açoitados.
\VS{23}E tendo sido lhes dado muitos açoites, os lançaram na prisão, mandando ao carcereiro que os guardasse em segurança.
\VS{24}O qual, tendo recebido tal ordem, lançou-os na cela mais interna, e prendeu-lhes os pés no tronco.
\VS{25}E perto da meia-noite, Paulo e Silas
{\ADD{estavam}} orando, e cantando hinos a Deus;
{\ADD{e}} os
{\ADD{outros}} presos os escutavam.
\VS{26}E de repente houve um terremoto tão grande que os alicerces da prisão se moviam; e logo todas as portas se abriram, e todas as correntes
{\ADD{que prendiam}} a todos se soltaram.
\VS{27}E o carcereiro, tendo acordado e visto abertas todas as portas da prisão; puxou a espada,
{\ADD{e}} estava a ponto de se matar, pensando que os presos tinham fugido.
\VS{28}Mas Paulo clamou em alta voz, dizendo: Não te faças nenhum mal, porque todos nós estamos aqui.
\VS{29}E tendo pedido luzes, saltou para dentro, e termendo muito, ele se prostrou diante de Paulo e Silas.
\VS{30}E levando-os para fora, disse: Senhores, o que me é necessário fazer para eu me salvar?
\VS{31}E eles
{\ADD{lhe}} disseram: Crê no Senhor Jesus Cristo, e serás salvo, tu e a tua casa.
\VS{32}E lhe falaram da palavra do Senhor, e a todos os que estavam em sua casa.
\VS{33}E ele, tomando-os consigo, naquela mesma hora da noite, lavou
{\ADD{-lhes}} as feridas dos açoites, e logo foi batizado, ele e todos os seus.
\VS{34}E tendo os levado a sua casa, pôs
{\ADD{comida}} diante deles à mesa; e alegrou-se muito, tendo crido em Deus com toda a sua casa.
\VS{35}E sendo
{\ADD{já}} de dia, os magistrados mandaram aos guardas, dizendo: Solta aqueles homens.
\VS{36}E o carcereiro anunciou estas palavras a Paulo,
{\ADD{dizendo}} : Os magistrados têm mandado vos soltar; portanto agora saí, e ide em paz.
\VS{37}Mas Paulo lhes disse: Eles nos açoitaram publicamente, e sem sermos sentenciados, sendo nós homens romanos, lançaram-nos na prisão, e agora nos lançam fora às escondidas?
{\ADD{Assim}} não! Mas que eles mesmos venham e nos tirem.
\VS{38}E os guardas voltaram para dizer aos magistrados estas palavras; e eles temeram ao ouvirem que eram romanos.
\VS{39}E tendo vindo, rogaram-lhes; e tirando-os, pediram
{\ADD{-lhes}} que saíssem da cidade.
\VS{40}E eles, tendo saído da prisão, entraram
{\ADD{na casa}} de Lídia; e vendo aos irmãos, consolaram-lhes; e saíram.

\par }\Chap{17}{\PP \VerseOne{1}E viajando por Anfípolis e Apolônia, vieram a Tessalônica, onde havia uma sinagoga de judeus.
\VS{2}E Paulo, como era de
{\ADD{seu}} costume, entrou a eles, e por três sábados argumentava com eles pelas Escrituras.
\VS{3}Declarando
{\ADD{-as}} ,e propondo
{\ADD{-lhes}} ,que era necessário que o Cristo morresse, e ressuscitasse dos mortos; e que (
{\ADD{dizia ele}} ) este Jesus é o Cristo, a quem eu vos anuncio.
\VS{4}E alguns deles creram, e se ajuntaram a Paulo e Silas; e dos gregos devotos grande multidão; e não poucas das mulheres principais.
\VS{5}Mas os judeus incrédulos, movidos de inveja, tomaram consigo alguns homens malignos dos mercados, e juntando uma multidão, tumultuaram a cidade; e atacando a casa de Jasão, procuravam trazê-los ao povo.
\VS{6}E não os achando, puxaram a Jasão, e a alguns irmãos às maiores autoridades da cidade, clamando: Estes que tem perturbado ao mundo também vieram até aqui.
\VS{7}Aos quais Jasão tem recolhido, e todos estes fazem contra as ordens de César, dizendo que há outro rei,
{\ADD{chamado}} Jesus.
\VS{8}E eles tumultuaram a multidão, e às autoridades da cidade, que ouviam estas coisas.
\VS{9}Mas tendo recebido fiança de Jasão e dos demais, eles os soltaram.
\VS{10}E logo os irmãos enviaram de noite a Paulo e a Silas até Bereia; os quais, ao chegarem lá, foram à sinagoga dos judeus.
\VS{11}E estes foram mais nobres do que os que estavam em Tessalônica,
{\ADD{porque}} os quais receberam a palavra com toda boa vontade, examinando a cada dia as escrituras,
{\ADD{para ver}} se estas coisas eram assim.
\VS{12}Portanto muitos deles realmente creram, e das mulheres gregas honradas; e dos homens, não poucos.
\VS{13}Mas quando os judeus de Tessalônica souberam que também em Bereia a palavra de Deus era anunciada por Paulo, vieram também para lá, e incitaram as multidões.
\VS{14}Mas então no mesmo instante os irmãos despediram a Paulo, para que fosse ao mar; mas Silas e Timóteo continuaram ali.
\VS{15}E os que conduziram a Paulo, levaram-no até Atenas; e tendo recebido ordem para Silas e Timóteo, para que viessem a ele o mais rápido
{\ADD{que pudessem}} ,eles foram embora.
\VS{16}E enquanto Paulo os esperava em Atenas, seu espírito se incomodava dentro dele, ao ver a cidade tão dedicada à idolatria.
\VS{17}Então ele disputava muito na sinagoga, com os judeus, e com os religiosos; e no mercado a cada dia, com os que vinham
{\ADD{até ele}} .
\VS{18}E alguns dos filósofos epicureus e estoicos discutiam com ele; e uns diziam: O que quer dizer este tagarela?E outros: Parece que ele é pregador de deuses estranhos.Porque ele lhes anunciava o Evangelho de Jesus e a ressurreição.
\VS{19}E tomando-o, trouxeram
{\ADD{-no}} ao areópago, dizendo: Podemos nós saber que doutrina nova é esta que tu falas?
\VS{20}Porque tu trazes coisas estranhas aos nossos ouvidos; então queremos saber o que isto quer dizer.
\VS{21}(E todos os atenienses e visitantes estrangeiros não se ocupavam de nenhuma outra coisa, a não ser em dizer e ouvir alguma novidade).
\VS{22}E Paulo, estando no meio do areópago, disse: Homens atenienses, eu vejo em tudo como vós sois muito religiosos;
\VS{23}Porque enquanto eu passava
{\ADD{pela cidade}} e via vossos santuários, achei também um altar, em que estava escrito: AO DEUS DESCONHECIDO. Este a quem vós prestais devoção sem conhecer,
{\ADD{este é o que}} eu vos anuncio;
\VS{24}O Deus que fez o mundo, e todas as coisas que nele
{\ADD{há}} ; este, sendo Senhor do céu e da terra, não habita em tempos feitos por mãos.
\VS{25}E nem também é servido por mãos humanas;
{\ADD{como que}} necessitasse de alguma coisa;
{\ADD{porque}} ele
{\ADD{é que}} dá a todos vida, respiração, e todas as coisas;
\VS{26}E de um sangue ele fez toda nação humana, para habitarem sobre a face da terra, determinando os tempos desde antes ordenados, e o limites da morada
{\ADD{deles}} ;
\VS{27}Para que buscassem ao Senhor, se talvez pudessem apalpá-lo e encontrá-lo,apesar dele não estar longe de cada um de nós.
\VS{28}Porque nele vivemos, e nos movemos, e somos; assim como também alguns de vossos poetas disseram; porque também nós somos descendência dele.
\VS{29}Sendo então descendência de Deus, nós não devemos pensar que a divindade seja semelhante a ouro, ou a prata, ou a pedra esculpida por artifício e imaginação humana.
\VS{30}Portanto Deus, tendo desconsiderado os tempos da
{\ADD{vossa}} ignorância, agora anuncia a todos as pessoas, em todo lugar, para que se arrependam.
\VS{31}Porque ele tem estabelecido um dia em que ele julgará ao mundo em justiça por meio do homem a quem determinou; dando certeza a todos, tendo o ressuscitado dos mortos.
\VS{32}E ao ouvirem da ressurreição dos mortos, alguns zombavam; e outros diziam: Nós ouviremos sobre isto de ti na próxima vez.
\VS{33}E assim Paulo saiu do meio deles.
\VS{34}Porém, tendo chegado alguns homens até ele, creram; entre os quais estava também Dionísio o areopagita, e uma mulher de nome Dâmaris, e outros com eles.

\par }\Chap{18}{\PP \VerseOne{1}E depois disto ele partiu de Atenas, e veio a Corinto.
\VS{2}E achando a um certo judeu, de nome Áquila, natural de Ponto, que recentemente tinha vindo da Itália, e Priscila, sua mulher (porque Cláudio tinha mandado que todos os judeus saíssem de Roma), veio até eles.
\VS{3}E porque era do mesmo ofício, ficou com eles, e trabalhava; porque tinham o ofício de fazerem tendas.
\VS{4}E ele disputava na sinagoga a cada sábado; e persuadia a judeus e a gregos.
\VS{5}E quando Silas e Timóteo desceram da Macedônia, Paulo foi pressionado pelo Espírito, dando testemunho aos judeus
{\ADD{de que}} o Cristo
{\ADD{era}} Jesus.
\VS{6}Mas tendo eles resistido e blasfemado, ele sacudiu as roupas, e lhes disse: Vosso sangue
{\ADD{seja}} sobre vossa cabeça; eu estou limpo; e a partir de agora irei aos gentios.
\VS{7}E tendo saído dali, entrou na cada de um, de nome Justo, que servia a Deus, cuja casa era vizinha à sinagoga.
\VS{8}E Crispo, o chefe da sinagoga, creu no Senhor com toda a sua casa; e muitos dos coríntios, tendo ouvido, creram e foram batizados.
\VS{9}E o Senhor disse em visão de noite a Paulo: Não temas, mas fala, e não te cales.
\VS{10}Porque eu estou contigo, e ninguém porá
{\ADD{mão}} em ti para te fazer mal, porque eu tenho muito povo nesta cidade.
\VS{11}E ele ficou
{\ADD{ali por}} um ano e seis meses, ensinando entre eles a palavra de Deus.
\VS{12}Mas sendo Gálio o procônsul da Acaia, os judeus se levantaram em concordância contra Paulo, e o trouxeram ao tribunal,
\VS{13}Dizendo: Este persuade as pessoas a servirem a Deus contra a Lei.
\VS{14}E Paulo, querendo abrir a boca, Gálio disse aos judeus: Se houvesse algum mau ato ou crime grande, ó judeus, com razão eu vos suportaria;
\VS{15}Mas se a questão é de palavra
{\ADD{s}} , e de nomes, e da Lei que há entre vós, vede
{\ADD{-o}} vós mesmos; porque destas coisas eu não quero ser juiz.
\VS{16}E ele os tirou do tribunal.
\VS{17}Mas todos os gregos, tomando a Sóstenes, o chefe da sinagoga, feriram
{\ADD{-no}} diante do tribunal; e a nada destas coisas Gálio dava importância.
\VS{18}E Paulo, ficando ali ainda muitos dias, ele se despediu dos irmãos, e dali navegou para a Síria, e juntos com ele
{\ADD{estavam}} Priscila e Áquila; tendo rapado a cabeça em Cencreia, porque ele tinha
{\ADD{feito}} voto.
\VS{19}E chegou a Éfeso, e os deixou ali; mas ele, entrando na sinagoga, disputava com os judeus.
\VS{20}E eles, pedindo{\ADD{-lhe}} que continuasse com eles por mais
{\ADD{algum}} tempo, ele não concordou.
\VS{21}Porém despediu-se deles, dizendo: De toda maneira tenho que estar na festa que vem em Jerusalém; mas outra vez, se Deus quiser, voltarei a vós.E ele saiu de Éfeso.
\VS{22}E tendo vindo a Cesareia, subiu e, saudando à igreja, desceu a Antioquia.
\VS{23}E passando
{\ADD{ali}} algum tempo, ele partiu, passando em sequência pela região da Galácia e Frígia, firmando a todos os discípulos.
\VS{24}E chegou a Éfeso um certo judeu, de nome Apolo, natural de Alexandria, homem eloquente, bem capacitado nas Escrituras.
\VS{25}Este era instruído no caminho do Senhor; e fervoroso de espírito, falava e ensinava corretamente as coisas do Senhor;
{\ADD{ainda que}} soubesse somente o batismo de João.
\VS{26}E este começou a falar ousadamente na sinagoga; e Áquila e Priscila, ao o ouvirem, tomaram-no consigo, e explicaram mais detalhadamente o caminho de Deus.
\VS{27}E ele, querendo passar a Acaia, os irmãos o exortaram,
{\ADD{e}} escreveram aos discípulos que o recebessem; o qual, tendo chegado, auxiliou muito aos que tinham crido pela graça.
\VS{28}Porque vigorosamente ele provava publicamente os erros dos judeus, mostrando pelas Escrituras que Jesus era o Cristo.

\par }\Chap{19}{\PP \VerseOne{1}E enquanto Apolo estava em Corinto, aconteceu que, tendo Paulo passado por todas as regiões altas, ele veio a Éfeso; e achando
{\ADD{ali}} alguns discípulos,
\VS{2}Disse-lhes: Vós
{\ADD{já}} recebestes o Espírito Santo
{\ADD{desde que}} crestes? E eles lhe disseram: Nós nem tínhamos ouvido falar que havia Espírito Santo.
\VS{3}E ele lhes disse: Em que vós fostes batizados? E eles disseram: No batismo de João.
\VS{4}Então Paulo disse: João verdadeiramente batizou
{\ADD{com}} o batismo do arrependimento, dizendo ao povo que cressem naquele que viria após ele, isto é, em Jesus Cristo.
\VS{5}Ao ouvirem
{\ADD{isto}} , eles foram batizados no nome do Senhor Jesus.
\VS{6}E Paulo, impondo-lhes as mãos, veio sobre eles o Espírito Santo; e falavam em línguas, e profetizavam.
\VS{7}E todos os homens eram cerca de doze.
\VS{8}E ele, entrando na sinagoga, falava ousadamente durante três meses, disputando e persuadindo as coisas do Reino de Deus.
\VS{9}Mas quando alguns se endureceram, e não creram, e falando mal do Caminho
\FTNT{*}{{\FR 19:9 }
Isto é, a fé em Jesus
} diante da multidão, ele se desviou deles; e separou aos discípulos, disputando a cada dia na escola de um certo Tirano.
\VS{10}E isto aconteceu durante dois anos; de tal maneira que todos os que habitavam na Ásia tinham ouvido a palavra do Senhor Jesus, tanto judeus como gregos.
\VS{11}E Deus fazia milagres extraordinários pelas mãos de Paulo,
\VS{12}De tal maneira que até os lenços e aventais de seu corpo eram levados aos enfermos, e as doenças os deixavam, e os espíritos malignos saíam deles.
\VS{13}E alguns exorcistas dos judeus, itinerantes, tentaram invocar o nome do Senhor Jesus sobre os que tinham espíritos malignos, dizendo: Nós vos repreendemos por Jesus, a quem Paulo prega.
\VS{14}E eram sete filhos de Ceva, judeu, chefe dos sacerdotes, os que faziam isto.
\VS{15}Mas o espírito maligno respondeu: Eu conheço Jesus, e sei quem é Paulo; mas vós, quem sois?
\VS{16}E o homem em quem estava o espírito maligno saltou sobre eles, e dominando-os, foi mais forte do que eles; de tal maneira que eles fugiram nus e feridos daquela casa.
\VS{17}E isto se fez conhecido a todos que habitavam em Éfeso, tanto a judeus como a gregos; e caiu temor sobre todos eles; e
{\ADD{assim}} foi engrandecido o nome do Senhor Jesus.
\VS{18}E muitos dos que criam vinham, confessando e declararando as suas atitudes.
\VS{19}Também muitos dos que praticavam ocultismo trouxeram seus livros, e
{\ADD{os}} queimaram na presença de todos; e calcularam o preço deles, e acharam que
{\ADD{custavam}} cinquenta mil
{\ADD{moedas}} de prata.
\VS{20}Assim a palavra do Senhor crescia e prevalecia poderosamente.
\VS{21}E quando se cumpriram estas coisas, Paulo propôs em espírito que, passando pela Macedônia e Acaia, ir até Jerusalém, dizendo: Depois de eu estar lá, também tenho que ver Roma.
\VS{22}E ele, tendo enviado à Macedônia dois daqueles que o serviam, Timóteo e Erasto, ele ficou por
{\ADD{algum}} tempo na Ásia.
\VS{23}Mas naquele tempo aconteceu um não pequeno alvoroço quanto ao Caminho
\FTNT{†}{{\FR 19:23 }
Isto é, a fé em Jesus
} .
\VS{24}Porque um certo artífice de prata, de nome Demétrio, que fazia objetos de prata para o templo de Diana
\FTNT{‡}{{\FR 19:24 }
Outro nome da mesma divindade era Ártemis
} , e dava não pouco lucro aos artesãos.
\VS{25}Aos quais, tendo os reunido com os trabalhadores de semelhantes coisas, disse: Homens, vós sabeis que deste ofício temos nossa prosperidade.
\VS{26}E vós estais vendo e ouvindo que este Paulo, não somente em Éfeso, mas até quase em toda a Ásia, tem persuadido e apartado uma grande multidão, dizendo que não são deuses os que são feitos com as mãos.
\VS{27}E não somente há perigo de que isto torne
{\ADD{nosso}} ofício em desprezo, mas também que
{\ADD{até}} o tempo da grande deusa Diana seja considerado inútil, e que a grandiosidade dela, a quem toda a Ásia e o mundo venera, venha a ser destruída.
\VS{28}E eles, ao ouvirem
{\ADD{estas coisas}} ,encheram-se de ira, e clamaram, dizendo: Grande é a Diana dos efésios!
\VS{29}E toda a cidade se encheu de tumulto, e em concordância eles invadiram ao teatro, tomando consigo Gaio e Aristarco, macedônios, companheiros de Paulo na viagem.
\VS{30}E Paulo, querendo comparecer diante do povo, os discípulos não o permitiram.
\VS{31}E também alguns dos líderes da Ásia, que eram amigos dele, enviaram-lhe
{\ADD{aviso}} , rogando{\ADD{-lhe}} para que não se apresentasse no teatro.
\VS{32}Então gritavam,
{\ADD{alguns de uma maneira}} , outros de outra; porque a aglomeração estava confusa; e a maioria não sabia por que causa estavam aglomerados.
\VS{33}E tiraram da multidão a Alexandre, os judeus o pondo para a frente. E Alexandre, acenando com a mão, queria
{\ADD{se}} defender ao povo.
\VS{34}Mas ao saberem que ele era judeu, levantou-se uma voz de todos, clamando por cerca de duas horas: Grande é a Diana dos efésios!
\VS{35}E o escrivão, tendo apaziguado a multidão, disse: Homens efésios, quem não sabe que a cidade dos efésios é a guardadora do templo da grande deusa Diana, e da
{\ADD{imagem}} que desceu do céu?
\VS{36}Portanto, sendo que estas coisas não podem ser contraditas, vós deveis vos acalmar, e nada façais precipitadamente;
\VS{37}Porque vós trouxestes
{\ADD{aqui}} estes homens, que nem são sacrílegos, nem blasfemam de vossa deusa.
\VS{38}Então se Demétrio e os artesãos que estão com ele tem algum assunto contra ele, os tribunais estão abertos, e há procônsules; que se acusem uns aos outros.
\VS{39}E se procurais alguma outra coisa, será decidido em uma reunião legalizada.
\VS{40}Porque corremos perigo de que hoje sejamos acusados de rebelião, tendo causa nenhuma para dar como explicação para este tumulto.
\VS{41}E tendo dito isto, despediu o ajuntamento.

\par }\Chap{20}{\PP \VerseOne{1}E tendo acabado o tumulto, Paulo chamou a si os discípulos e, abraçando-os, saiu para ir à Macedônia.
\VS{2}E tendo passado por aquelas regiões, e exortando-os com muitas palavras, ele veio à Grécia.
\VS{3}E ficando
{\ADD{ali}} por três meses, e havendo contra ele uma cilada posta pelos judeus quando ele estava a ponto de navegar para a Síria, ele decidiu voltar pela Macedônia.
\VS{4}E o acompanhou até a Ásia Sópater, de Bereia; e dos tessalonicenses, Aristarco e Segundo, e Gaio de Derbe, e Timóteo; e dos da Ásia, Tíquico e Trófimo.
\VS{5}Estes, indo adiante, nos esperaram em Trôade.
\VS{6}E depois dos dias dos
{\ADD{pães}} não fermentados, nós navegamos de Filipos, e em cinco dias viemos até eles, onde ficamos por sete dias.
\VS{7}E no primeiro
{\ADD{dia}} da semana, tendo os discípulos se reunido para partir o pão, Paulo discutia com eles, estando para partir no dia seguinte; e ele estendeu a discussão até a meia noite.
\VS{8}E havia muitas luminárias no compartimento onde estavam reunidos.
\VS{9}E estando um certo rapaz, de nome Êutico, sentado em uma janela, tendo sido tomado por um sono profundo,
{\ADD{e}} ,estando Paulo
{\ADD{ainda}} falando por muito
{\ADD{tempo}} ,{\ADD{Êutico}} ,derrubado pelo sono, caiu desde o terceiro andar abaixo; e foi levantado morto.
\VS{10}Mas Paulo, tendo descido, debruçou sobre ele e, abraçando
{\ADD{-o}} ,disse: Não fiqueis perturbados, porque sua alma
{\ADD{ainda}} está nele.
\VS{11}E
{\ADD{voltou}} a subir, e tendo partido e experimentado o pão, ele falou longamente até o nascer do dia; e assim ele partiu.
\VS{12}E trouxeram o rapaz vivo, e ficaram não pouco consolados.
\VS{13}E nós, tendo ido adiante ao navio, navegamos até Assôs, onde estaríamos para receber a Paulo, porque assim ele tinha ordenado; e ele ia a pé.
\VS{14}E quando ele se encontrou conosco em Assôs, nós o tomamos, e fomos a Mitilene.
\VS{15}E navegando dali, chegamos no
{\ADD{dia}} seguinte em frente a Quios; e no outro dia aportamos em Samos; e ficando em Trogílio, no dia seguinte viemos a Mileto.
\VS{16}Porque Paulo tinha decidido navegar
{\ADD{desviando-se}} de Éfeso, para não lhe haver de gastar tempo na Ásia; porque ele se apressava para estar em Jerusalém do dia de Pentecostes, caso lhe fosse possível.
\VS{17}Mas ele enviou
{\ADD{mensagem}} desde Mileto até Éfeso, chamando aos anciãos
\FTNT{*}{{\FR 20:17 }
anciãos = equiv. presbíteros
} da igreja.
\VS{18}E quando vieram a
{\ADD{Paulo}} ,ele lhes disse: Vós sabeis que desde o primeiro dia que entrei na Ásia,
{\ADD{o modo}} como eu estive todo
{\ADD{aquele}} tempo convosco;
\VS{19}Servindo ao Senhor com toda humildade, muitas lágrimas, e tentações, que sobrevieram a mim pelas ciladas dos judeus;
\VS{20}Como eu, daquilo que
{\ADD{vos}} era proveitoso, nada deixei de anunciar a vós, e ensinar publicamente e pelas casas;
\VS{21}Dando testemunho, tanto a judeus como a gregos, do arrependimento para
{\ADD{se converter}} a Deus, e da fé em nosso Senhor Jesus Cristo.
\VS{22}E agora eis que, estando eu atado ao Espírito, estou indo a Jerusalém, não sabendo o que me acontecerá;
\VS{23}A não ser pelo que o Espírito Santo em cada cidade
{\ADD{me}} dá testemunho, dizendo que prisões e aflições me esperam.
\VS{24}Mas de nenhuma
{\ADD{dessas}} coisas eu dou importância, nem tenho minha vida por preciosa, para que com alegria eu cumpra minha carreira, e o trabalho que recebi do Senhor Jesus, para dar testemunho do Evangelho da graça de Deus.
\VS{25}E agora, eis que eu sei que todos vós, a quem eu passei pregando o Reino de Deus, não vereis mais o meu rosto.
\VS{26}Portanto eu vos dou claro testemunho de que eu estou limpo do sangue de todos
{\ADD{vós}} ;
\VS{27}Porque eu não deixei de vos anunciar todo o conselho de Deus;
\VS{28}Portanto prestai atenção por vós mesmos, e por todo o rebanho sobre os quais o Espírito Santo tem vos posto como supervisores,
\FTNT{†}{{\FR 20:28 }
supervisores = equiv. bispos
} para apascentardes a igreja de Deus, a qual ele adquiriu por meio de seu próprio sangue.
\VS{29}Porque isto eu sei, que depois de minha partida, entrarão entre vós lobos cruéis, que não pouparão ao rebanho;
\VS{30}E que dentre vós mesmos se levantarão homens a falarem coisas perversas, para atraírem após si aos discípulos.
\VS{31}Por isso vigiai, lembrando que por três anos, noite e dia eu não parei de vos alertar com lágrimas a cada um
{\ADD{de vós}} .
\VS{32}E agora, irmãos, eu vos entrego a Deus, e à palavra de sua graça; ele que é poderoso para vos edificar e vos dar herança entre todos os santificados.
\VS{33}Eu não cobicei de ninguém a prata, nem ouro, nem roupa.
\VS{34}E vós mesmos sabeis que, para as minhas necessidades e as dos que estavam comigo, estas
{\ADD{minhas}} mãos me serviram.
\VS{35}Em tudo eu vos tenho mostrado que trabalhando assim, é necessário dar suporte aos enfermos; e lembrar-se das palavras do Senhor Jesus, que disse: Mais bem-aventurado é dar do que receber.
\VS{36}E tendo dito isto, pondo-se de joelhos, ele orou com todos eles.
\VS{37}E houve um grande pranto de todos; e reclinando-se sobre o pescoço de Paulo, beijavam-no;
\VS{38}Entristecendo-se muito, principalmente pela palavra que ele tinha dito, que não mais veriam o rosto dele; e o acompanharam até o navio.

\par }\Chap{21}{\PP \VerseOne{1}E quando aconteceu de termos saído deles, e navegado, percorremos diretamente, e viemos a Cós, e
{\ADD{n}} o
{\ADD{dia}} seguinte a Rodes, e dali a Pátara.
\VS{2}E tendo achado um navio que passava para a Fenícia, nós embarcamos nele, e partimos.
\VS{3}E tendo Chipre à vista, e deixando-a à esquerda, navegamos para a Síria, e viemos a Tiro; porque o navio tinha de deixar ali sua carga.
\VS{4}E nós ficamos ali por sete dias; e achamos aos discípulos, os quais diziam pelo Espírito a Paulo, que não subisse a Jerusalém.
\VS{5}E tendo passado
{\ADD{ali}} aqueles dias, nós saímos e seguimos nosso caminho, acompanhando-nos todos com
{\ADD{suas}} mulheres e filhos até fora da cidade; e postos de joelhos na praia, oramos.
\VS{6}E saudando-nos uns aos outros, subimos ao navio; e eles voltaram para as
{\ADD{casas}} deles.
\VS{7}E nós, acabada a navegação de Tiro, chegamos a Ptolemaida; e tendo saudado aos irmãos, ficamos com eles por um dia.
\VS{8}E
{\ADD{n}} o
{\ADD{dia}} seguinte, Paulo e nós que estávamos com ele, saindo dali, viemos a Cesareia; e entrando na casa de Filipe, o evangelista (que era
{\ADD{um}} dos sete), nós ficamos com ele.
\VS{9}E este tinha quatro filhas, que profetizavam.
\VS{10}E ficando nós
{\ADD{ali}} por muitos dias, desceu da Judeia um profeta, de nome Ágabo;
\VS{11}E ele, tendo vindo a nós, e tomando a cinta de Paulo, e atando-se os pés e as mãos, disse: Isto diz o Espírito Santo: Assim os judeus em Jerusalém atarão ao homem a quem
{\ADD{pertence}} esta cinta, e o entregarão nas mãos dos gentios.
\VS{12}E nós, tendo ouvido isto, rogamos a ele, tanto nós como os que eram daquele lugar, que ele não subisse a Jerusalém.
\VS{13}Mas Paulo respondeu: O que vós estais fazendo, ao chorarem e afligirem o meu coração? Porque eu estou pronto, não somente para ser atado, mas até mesmo para morrer em Jerusalém, pelo nome do Senhor Jesus.
\VS{14}E como ele não deixou ser persuadido, nós nos aquietamos, dizendo: Faça-se a vontade do Senhor.
\VS{15}E depois daqueles dias, nós nos arrumamos, e subimos a Jerusalém.
\VS{16}E foram também conosco
{\ADD{alguns}} dos discípulos de Cesareia, trazendo
{\ADD{consigo}} a um certo Mnáson, cipriota, discípulo antigo, com o qual íamos nos hospedar.
\VS{17}E quando chegamos a Jerusalém, os irmãos nos receberam com muito boa vontade.
\VS{18}E
{\ADD{n}} o
{\ADD{dia}} seguinte, Paulo entrou conosco a
{\ADD{casa de}} Tiago, e todos os anciãos vieram ali.
\VS{19}E tendo os saudado, ele
{\ADD{lhes}} contou em detalhes o que Deus tinha feito entre os gentios por meio do trabalho dele.
\VS{20}E eles, ao ouvirem, glorificaram ao Senhor, e lhe disseram: Tu vês, irmão, quantos milhares de judeus há que creem, e todos são zelosos da lei.
\VS{21}E foram informados quanto a ti, que a todos os judeus, que estão entre os gentios,
{\ADD{que}} tu ensinas a se afastarem de Moisés, dizendo que não devem circuncidar
{\ADD{seus}} filhos, nem andar segundo os costumes.
\VS{22}Então o que se fará? Em todo caso a multidão deve se ajuntar, porque ouvirão que tu
{\ADD{já}} chegaste.
\VS{23}Portanto faça isto que te dizemos: temos quatro homens que fizeram voto.
\VS{24}Toma contigo a estes, e purifica-te com eles, e paga os gastos deles, para que rapem a cabeça, e todos saibam que não há nada do que foram informados sobre ti, mas sim,
{\ADD{que}} tu mesmo andas guardando a Lei.
\VS{25}E quanto aos que creem dentre os gentios, nós já escrevemos, julgando que nada disto guardassem, a não ser somente que se abstenham do que se abstenham das coisas sacrificadas aos ídolos, e do sangue, e da
{\ADD{carne}} sufocada, e do pecado sexual.
\VS{26}Então Paulo, tendo tomado consigo a aqueles homens, e purificando-se com eles no dia seguinte, entrou no Templo, anunciando que
{\ADD{já}} estavam cumpridos dos dias da purificação, quando fosse oferecida por eles, cada um
{\ADD{sua}} oferta.
\VS{27}E tendo os sete dias quase quase completados, os judeus da Ásia, vendo-o, tumultuaram a todo o povo, e lançaram as mãos sobre ele
{\ADD{para o deterem}} ,
\VS{28}Clamando: Homens israelitas, ajudai
{\ADD{-nos}} ; este é o homem que por todos os lugares ensina a todos contra o
{\ADD{nosso}} povo, e contra a Lei, e
{\ADD{contra}} este lugar; e além disto, ele também pôs gregos dentro do Templo, e contaminou este santo lugar!
\VS{29}(Porque antes eles tinham visto na cidade a Trófimo junto dele, ao qual pensavam que Paulo tinha trazido para dentro do Templo).
\VS{30}E toda a cidade se tumultuou, e houve um ajuntamento do povo; e tendo detido a Paulo, trouxeram-no para fora do Templo; e logo as portas foram fechadas.
\VS{31}E eles, procurando matá-lo, a notícia chegou ao comandante e aos soldados, de que toda Jerusalém estava em confusão.
\VS{32}O qual, tendo tomado logo consigo soldados e centuriões, correu até eles. E eles, vendo ao comandante e aos soldados, pararam de ferir a Paulo.
\VS{33}Então o comandante, tendo se aproximado, prendeu-o, e mandou que ele fosse atado em duas correntes; e perguntou quem ele era, e o que ele tinha feito.
\VS{34}E na multidão clamavam
{\ADD{uns de uma maneira}} ,e outros de outra maneira; mas
{\ADD{como}} ele não podia saber com certeza por causa do tumulto, ele mandou que o levassem para a área fortificada.
\VS{35}E ele, tendo chegado às escadas, aconteceu que ele foi carregado pelos soldados, por causa da violência da multidão.
\VS{36}Porque a multidão do povo
{\ADD{o}} seguia, gritando: Tragam-no para fora!
\VS{37}E Paulo, estando perto de entrar na área fortificada, disse ao comandante: É permitido a mim te falar alguma coisa?E ele disse: Tu sabes grego?
\VS{38}Por acaso não és tu aquele egípcio, que antes destes dias tinha levantando uma rebelião, e levou ao deserto quatro mil homens assassinos?
\VS{39}Mas Paulo lhe disse: Na verdade eu sou um homem judeu de Tarso, cidade não pouca importância da Cilícia; mas eu te rogo para que tu me permitas falar ao povo.
\VS{40}E tendo
{\ADD{lhe}} permitido, Paulo pôs-se de pé nas escadas, acenou com a mão ao povo; e tendo havido grande silêncio, falou
{\ADD{-lhes}} em língua hebraica, dizendo:

\par }\Chap{22}{\PP \VerseOne{1}Homens irmãos, e pais, ouvi agora minha defesa para convosco.
\VS{2}E tendo ouvido que ele lhes falava em língua hebraica, fizeram ainda mais silêncio. E ele disse:
\VS{3}Eu verdadeiramente sou um homem judeu, nascido em Tarso da Cilícia, e criado nesta cidade aos pés de Gamaliel, ensinado ao mais correto modo da Lei paterna
{\ADD{e}} zeloso de Deus, assim como todos vós sois hoje.
\VS{4}{\ADD{Eu}} ,que persegui este caminho até a morte, atando tanto a homens como a mulheres, e
{\ADD{os}} entregando a prisões.
\VS{5}Assim como também o sumo sacerdote me é testemunha, e todo o conselho dos anciãos; dos quais eu, tendo tomado cartas para os irmãos, fui a Damasco para que os que estivessem ali, eu também os trouxesse amarrados a Jerusalém, para que fossem castigados.
\VS{6}Mas aconteceu que, estando eu no caminho, e chegando perto de Damasco, quase ao meio dia, de repente uma grande luz do céu brilhou ao redor de mim.
\VS{7}Eu cai ao chão, e ouvi uma voz, que me dizia: Saulo, Saulo, por que me persegues?
\VS{8}E eu respondi: Quem és, Senhor? E ele me disse: Eu sou Jesus, o nazareno, a quem tu persegues.
\VS{9}E os que estavam comigo verdadeiramente viram a luz, e ficaram muito atemorizados; mas eles não ouviram a voz daquele que falava comigo.
\VS{10}E eu disse: Que farei, Senhor? E o Senhor me disse: Levanta-te, e vai a Damasco, e ali te será dito tudo o que te é ordenado fazer.
\VS{11}E quando eu não
{\ADD{conseguia}} ver, por causa da glória daquela luz, eu fui levado pela mão dos que estavam comigo, e
{\ADD{assim}} cheguei a Damasco.
\VS{12}E um certo Ananias, homem devoto conforme a Lei, que tinha
{\ADD{bom}} testemunho de todos os judeus que moravam
{\ADD{ali}} ;
\VS{13}Tendo vindo até mim, e ficando
{\ADD{diante de mim}} ,ele me disse: Irmão Saulo, recupere a vista; E naquela mesma hora eu
{\ADD{pude}} vê-lo.
\VS{14}E ele disse: O Deus de nossos pais te predeterminou para que tu conheças a vontade dele, e vejas aquele justo, e tu ouças a voz de sua boca.
\VS{15}Porque tu serás testemunha dele para com todos as pessoas, daquilo que tens visto e ouvido.
\VS{16}E agora, por que estás parado? Levanta-te, e sê batizado, e lava teus pecados, invocando o nome do Senhor.
\VS{17}E aconteceu a mim, tendo eu voltado a Jerusalém, e estando orando no Templo, veio-me um êxtase;
\VS{18}E eu vi aquele que me dizia: Apressa-te, e sai logo de Jerusalém, porque não aceitarão teu testemunho sobre mim.
\VS{19}E eu disse: Senhor, eles sabem que eu prendia e açoitava nas sinagogas aqueles que criam em ti.
\VS{20}E quando o sangue de Estêvão, tua testemunha, se derramava, eu também estava presente, e consentia em sua morte, e guardava as roupas daqueles que o matavam.
\VS{21}E ele me disse: Vai, porque eu te enviarei para longe, aos gentios.
\VS{22}E eles o ouviram até esta palavra, e
{\ADD{em seguida}} levantaram suas vozes, dizendo: Extermina-o da terra! Porque não é bom que ele viva.
\VS{23}E enquanto eles gritavam, tiravam
{\ADD{suas}} capas, e lançavm pó ao ar,
\VS{24}O comandante mandou que o levassem à área fortificada, dizendo que o interrogassem com açoites, para saber por que causa clamavam assim contra ele.
\VS{25}E quando estavam o atando com correias, Paulo disse ao centurião que estava ali: É lícito para vós açoitar a um homem romano, sem
{\ADD{ter sido}} condenado?
\VS{26}E o centurião, tendo ouvido
{\ADD{isto}} ,foi e avisou ao comandante, dizendo: Olha o que estás a ponto de fazer, porque este homem é romano.
\VS{27}E o comandante, tendo se aproximado, disse-lhe: Dize-me, tu és romano?E ele disse: Sim.
\VS{28}E o comandante respondeu: Eu com muita soma
{\ADD{de dinheiro}} obtive esta cidadania
{\ADD{romana}} .E Paulo disse: E eu
{\ADD{a tenho}} desde que nasci.
\VS{29}Então logo se afastaram dele aqueles que estavam para interrogá-lo; e até o comandante teve temor, ao entender que
{\ADD{Paulo}} era romano, e que tinha o atado.
\VS{30}E no{\ADD{dia}} seguinte, querendo saber corretamente a causa de por que ele era acusado pelos judeus, ele o soltou das correntes, e mandou vir aos chefes dos sacerdotes e todo o supremo conselho
\FTNT{*}{{\FR 22:30 }
supremo conselho
lit. sinédrio – o mais importante conselho ou tribunal para os judeus
} deles; e tendo trazido a Paulo, apresentou
{\ADD{-o}} diante deles.

\par }\Chap{23}{\PP \VerseOne{1}E Paulo, olhando fixamente para o supremo conselho,
\FTNT{*}{{\FR 23:1 }
supremo conselho
lit. sinédrio – o mais importante conselho ou tribunal para os judeus
} disse: Homens irmãos, com toda boa consciência eu tenho andado diante de Deus até o dia de hoje.
\VS{2}Mas o sumo sacerdote Ananias mandou aos que estavam perto dele, que o espancassem na boca.
\VS{3}Então Paulo lhe disse: Deus vai te espancar, parede caiada! Estás tu
{\ADD{aqui}} sentado para me julgar conforme a Lei, e contra a Lei mandas me espancarem?
\VS{4}E os que estavam ali disseram: Tu insultas ao sumo sacerdote de Deus?
\VS{5}E Paulo disse: Eu não sabia, irmãos, que ele era o sumo sacerdote; porque está escrito: Não falarás mal do chefe do teu povo.
\VS{6}E Paulo, sabendo que uma parte era de saduceus, e outra de fariseus, ele clamou no supremo conselho: Homens irmãos, eu sou fariseu, filho de fariseu; pela esperança e ressurreição dos mortos eu estou sendo julgado.
\VS{7}E ele, tendo dito isto, houve uma confusão entre os fariseus e saduceus; e a multidão se dividiu;
\VS{8}Porque os saduceus dizem que não há ressurreição, nem anjo ou espírito; mas os fariseus declaram ambas.
\VS{9}E houve uma grande gritaria; e levantando-se os escribas da parte dos fariseus, disputavam, dizendo: Nenhum mal achamos neste homem; e se algum espírito ou anjo falou com ele, não briguemos contra Deus.
\VS{10}E havendo grande confusão, o comandante, temendo que Paulo não fosse despedaçado por eles, mandou descer a tropa, e tirá-lo do meio deles, e levá-lo à área fortificada.
\VS{11}E
{\ADD{n}} a noite seguinte o Senhor, aparecendo-lhe, disse: Tem bom ânimo, Paulo! Porque assim como deste testemunho de mim em Jerusalém, assim é necessário que tu dês testemunho também em Roma.
\VS{12}E tendo vindo o dia, alguns dos judeus fizeram uma conspiração, e prestaram juramento sob pena de maldição, dizendo que não comeriam nem beberiam enquanto não matassem a Paulo.
\VS{13}E eram mais de quarenta os que fizeram este juramento.
\VS{14}Os quais foram até os chefes dos sacerdotes e os anciãos,
{\ADD{e}} disseram: Fizemos juramento sob pena de maldição, de que nada experimentaremos enquanto não matarmos a Paulo.
\VS{15}Agora, pois, vós, com o supremo conselho, informai ao comandante que amanhã ele o traga perante vós, como se fosse para que investigueis mais detalhadamente; e antes que ele chegue, estaremos prontos para o matar.
\VS{16}E o filho da irmã de Paulo, tendo ouvido esta cilada, veio e entrou na área fortificada, e avisou a Paulo.
\VS{17}E Paulo, tendo chamado a si um dos centuriões, disse: Leva este rapaz ao comandante, porque ele tem algo para lhe avisar.
\VS{18}Então ele o tomou, levou ao comandante, e disse: O prisioneiro Paulo, tendo me chamado, rogou
{\ADD{-me}} que eu te trouxesse este rapaz, que tem algo a te dizer.
\VS{19}E o comandante, tomando-o pela mão, e indo para um lugar reservado, perguntou
{\ADD{-lhe}} : O que tens para me avisar?
\VS{20}E ele disse: Os judeus combinaram de te pedirem que amanhã tu leves a Paulo ao supremo conselho, como se fosse para que lhe perguntem mais detalhadamente;
\VS{21}Porém tu, não acredites neles; porque mais de quarenta homens deles estão lhe preparando cilada, os quais sob pena de maldição fizeram juramento para não comerem nem beberem enquanto não o tiverem matado; e eles já estão preparados, esperando de ti a promessa.
\VS{22}Então o comandante despediu ao rapaz, mandando
{\ADD{-lhe}} : A ninguém digas que tu me revelaste estas coisas.
\VS{23}E ele, chamando a si certos dois dos centuriões, disse: Aprontai duzentos soldados para irem até Cesareia; e setenta cavaleiros, e duzentos arqueiros, a partir das terceira hora da noite.
\VS{24}E preparem animais para cavalgarem, para que pondo neles a Paulo, levem
{\ADD{-no}} a salvo ao governador Félix.
\VS{25}E ele
{\ADD{lhe}} escreveu uma carta, que continha este aspecto:
\VS{26}Cláudio Lísias, a Félix, excelentíssimo governador, saudações.
\VS{27}Este homem foi preso pelos judeus, e estando já a ponte de o matarem, eu vim com a tropa e
{\ADD{o}} tomei, ao ser informado que ele era romano.
\VS{28}E eu, querendo saber a causa por que o acusavam, levei-o ao supremo conselho deles.
\VS{29}O qual eu achei que acusavam de algumas questões da Lei deles; mas que nenhum crime digno de morte ou de prisão havia contra ele.
\VS{30}E tendo sido avisado de que os judeus estavam para pôr uma cilada contra este homem, logo eu
{\ADD{o}} enviei a ti, mandando também aos acusadores que diante de ti digam o que
{\ADD{tiverem}} contra ele. Que tu estejas bem.
\VS{31}Tendo então os soldados tomado a Paulo, assim como lhes tinha sido ordenado, trouxeram-no durante a noite a Antipátride.
\VS{32}E no
{\ADD{dia}} seguinte, deixando irem com ele os cavaleiros, voltaram à área fortificada.
\VS{33}Os quais, tendo chegado a Cesareia, e entregado a carta ao governador, apresentaram-lhe também a Paulo.
\VS{34}E o governador, tendo lido
{\ADD{a carta}} , perguntou de que província ele era; e ao entender que
{\ADD{era}} da Cilícia,
\VS{35}disse: “Eu te ouvirei quando também chegarem os teus acusadores”. E mandou que o guardassem no palácio
\FTNT{†}{{\FR 23:35 }
Ou: pretório
} de Herodes.

\par }\Chap{24}{\PP \VerseOne{1}E cinco dias depois, desceu o sumo sacerdote Ananias, com os anciãos e um certo orador Tértulo; os quais compareceram diante do governador contra Paulo.
\VS{2}E sendo chamado, Tértulo começou a acusá
{\ADD{-lo}} , dizendo:
\VS{3}{\ADD{Visto que}} há muita paz por causa de ti, e que por teu governo muitos bons serviços estão sendo feitos a esta nação, excelentíssimo Félix, totalmente e em todo lugar, com todo agradecimento o reconhecemos.
\VS{4}Mas para que eu não gaste muito o teu tempo, rogo
{\ADD{-te}} que tu nos ouças brevemente, conforme a tua clemência.
\VS{5}Porque nós temos achado que este homem
{\ADD{é}} uma peste, e levantador de rebeliões entre todos os judeus pelo mundo, e o principal defensor da seita dos nazarenos.
\VS{6}O qual também tentou profanar o Templo; ao qual também prendemos, e quisemos julgar conforme a nossa Lei.
\VS{7}Mas tendo vindo sobre
{\ADD{nós}} o comandante Lísias, com grande violência, tirou
{\ADD{-o}} das nossas mãos;
\VS{8}investigando-o tu mesmo, poderás entender todas estas coisas das quais o acusamos.
\VS{9}E os judeus também concordaram, dizendo serem estas coisas assim.
\VS{10}Mas Paulo, fazendo gesto ao governador para que falasse, respondeu: Sabendo eu que por muitos anos tu tens sido juiz desta nação, com maior ânimo eu me defendo.
\VS{11}Pois tu podes entender que há não mais que doze dias, eu tinha subido a Jerusalém para adorar.
\VS{12}E nem me acharam falando com alguém no Templo, nem incitando ao povo, nem nas sinagogas, nem na cidade.
\VS{13}E nem podem provar as coisas das quais agora estão me acusando.
\VS{14}Mas isto eu te confesso, que conforme o Caminho que eles chamam de seita, assim eu sirvo ao Deus dos
{\ADD{meus}} pais, crendo em tudo que está escrito na Lei e nos profetas;
\VS{15}Tendo esperança em Deus, ao qual estes mesmos também esperam, que vai haver ressurreição dos mortos, tanto dos justos como dos injustos.
\VS{16}E nisto eu pratico,
{\ADD{em que}} tanto para com Deus como para com os seres humanos eu sempre tenha uma consciência limpa.
\VS{17}E muitos anos depois, vim para fazer doações e ofertas à minha nação.
\VS{18}Nisto, tendo eu me purificado, nem com multidões, nem com tumulto, alguns judeus da Ásia me acharam,
\VS{19}Os quais era necessário que estivessem
{\ADD{aqui}} presentes diante de ti, se tivessem alguma coisa contra mim.
\VS{20}Ou digam estes mesmos, se eles acharam alguma má ação em mim quando eu estava perante o supremo conselho;
\VS{21}A não ser somente esta palavra com que, quando eu estava entre eles, clamei, que pela ressurreição dos mortos eu estou sendo julgado hoje por vós.
\VS{22}Tendo ouvido estas coisas, Félix, que sabia mais detalhadamente sobre o Caminho, adiou-lhes, dizendo: Quando o comandante Lísias descer, eu procurarei saber melhor de vossos assuntos.
\VS{23}E ele mandou ao centurião que guardassem a Paulo, e estivesse com
{\ADD{alguma}} liberdade, e impedir a ninguém dos seus
{\ADD{amigos}} lhe prestasse serviço, ou vir até ele.
\VS{24}E alguns dias depois, tendo vindo Félix com a mulher dele Drusila, que era judia, mandou chamar a Paulo, e o ouviu sobre a fé em Cristo.
\VS{25}E ele, tendo discursado sobre a justiça, o domínio próprio, e o julgamento que está por vir, Félix temeu,
{\ADD{e}} respondeu: Por agora vai; e tendo
{\ADD{outra}} oportunidade, eu te chamarei.
\VS{26}Ele também esperava que lhe fosse dado
{\ADD{algum}} dinheiro por Paulo, para que o soltasse; por isso ele também muitas vezes o mandava chamar, e conversava com ele.
\VS{27}Mas tendo completado dois anos, Félix teve por sucessor a Pórcio Festo. E Félix, querendo agradar dos judeus, deixou Paulo preso.

\par }\Chap{25}{\PP \VerseOne{1}Então Festo, tendo entrado na província, subiu dali três dias depois de Cesareia a Jerusalém.
\VS{2}E o sumo sacerdote e os líderes dos judeus compareceram diante dele contra Paulo, e lhe rogaram;
\VS{3}Pedindo favor contra ele, para que o fizesse vir a Jerusalém, preparando cilada para o matarem no caminho.
\VS{4}Mas Festo respondeu que Paulo estava guardado em Cesareia, e que ele logo estava indo
{\ADD{para lá}} .
\VS{5}Ele disse: Então aqueles dentre vós que podem, desçam com
{\ADD{igo}} , e se houver alguma coisa errada neste homem, acusem-no.
\VS{6}E ele, tendo ficado entre eles mais de dez dias, desceu a Cesareia; e tendo se sentado no tribunal no
{\ADD{dia}} seguinte, mandou que trouxessem a Paulo.
\VS{7}E tendo ele vindo, os judeus que haviam descido de Jerusalém
{\ADD{o}} rodearam, trazendo contra Paulo muitas e graves acusações, que não podiam provar.
\VS{8}Ele, disse em sua defesa: Eu não pequei nem contra a Lei dos judeus, nem contra o Templo, nem contra César, em coisa alguma.
\VS{9}Mas Festo, querendo agradar aos judeus, respondendo a Paulo, disse: Tu queres subir a Jerusalém e ser julgado sobre estas coisas diante de mim?
\VS{10}E Paulo disse: Eu estou diante do tribunal de César, onde eu tenho que ser julgado; a nenhum dos judeus eu fiz mal, assim como também tu sabes muito bem.
\VS{11}Porque se eu fiz algum mal, ou cometi algo digno de morte, eu não recuso morrer; mas se nada há das coisas que este me acusam, ninguém pode me entregar a eles. Eu apelo a César.
\VS{12}Então Paulo, tendo conversado com o Conselho, respondeu: Tu apelaste a César; a César irás.
\VS{13}E passados alguns dias, o Rei Agripa e Berenice vieram a Cesareia para saudar a Festo.
\VS{14}E quando tinham ficado ali muitos dias, Festo contou ao rei os assuntos de Paulo, dizendo: Um certo homem foi deixado
{\ADD{aqui}} preso por Félix;
\VS{15}Por causa do qual, estando eu em Jerusalém, os chefes dos sacerdotes e os anciãos dos judeus compareceram
{\ADD{a mim}} , pedindo julgamento contra ele.
\VS{16}Aos quais eu respondi não ser costume dos romanos entregar a algum homem à morte, antes que o acusado tenha seus acusadores face a face, e tenha oportunidade para se defender da acusação.
\VS{17}Portanto, tendo eles se reunido aqui, fazendo nenhum adiamento, no
{\ADD{dia}} seguinte, estando eu sentado no tribunal, mandei trazer ao homem.
\VS{18}Do qual os acusadores estando
{\ADD{aqui}} presentes, trouxeram como acusação nenhuma das coisas que eu suspeitava.
\VS{19}Mas tinham contra ele algumas questões relativas às próprias crenças deles, e de um certo morto Jesus, o qual Paulo afirmava estar vivo.
\VS{20}E eu, estando em duvida sobre
{\ADD{como}} interrogar esta causa, disse,
{\ADD{perguntando}} se ele queria ir a Jerusalém, e lá ser julgado sobre estas coisas.
\VS{21}Porém Paulo, tendo apelado para ser guardado ao interrogatório do imperador, mandei que o guardassem, até que eu o enviasse a César.
\VS{22}E Agripa disse a Festo: Eu também queria ouvir a este homem.E ele disse: E amanhã tu o ouvirás.
\VS{23}Então no dia seguinte, tendo vindo Agripa e Berenice, com muita pompa, e entrando no auditório com os comandantes e os homens mais importantes da cidade, trouxeram a Paulo por ordem de Festo.
\VS{24}E Festo disse: Rei Agripa, e todos os homens que estais presentes
{\ADD{aqui}} conosco, vós vedes este
{\ADD{homem}} , a quem toda a multidão dos judeus, tanto em Jerusalém como aqui, tem apelado a mim, clamando que ele não deve mais viver.
\VS{25}Mas tendo eu achado nada que ele tenha feito que fosse digno de morte, e também tendo ele mesmo apelado ao imperador, eu decidi enviá-lo.
\VS{26}Do qual eu não tenho coisa alguma certa para escrever ao
{\ADD{meu}} senhor; por isso que eu o trouxe diante de vós; e principalmente diante de ti, rei Agripa, para que, sendo feita a investigação, eu tenha algo para escrever.
\VS{27}Porque não me parece razoável enviar a um prisioneiro, sem também informar as acusações contra ele.

\par }\Chap{26}{\PP \VerseOne{1}E Agripa disse a Paulo: É permitido a ti falar por ti mesmo.Então Paulo, estendendo a mão, respondeu em sua defesa:
\VS{2}Eu me considero feliz, rei Agripa, de que diante de ti eu esteja hoje fazendo minha defesa de todas as coisas de que sou acusado pelos judeus;
\VS{3}Principalmente
{\ADD{por}} eu saber que tu sabes de todos os costumes e questões que há entre os judeus; por isso eu te rogo que tu me ouças com paciência.
\VS{4}Ora, a minha vida é conhecida por todos os judeus, desde a
{\ADD{minha}} juventude, que desde o princípio tem sido entre
{\ADD{os de}} minha nação em Jerusalém;
\VS{5}Eles me conhecem desde o começo, se quiserem testemunhar, de que conforme a mais rigorosa divisão de nossa religião, eu vivi
{\ADD{como}} fariseu.
\VS{6}E agora, pela esperança da promessa que por Deus foi dada aos nossos pais, eu estou
{\ADD{aqui}} sendo julgado.
\VS{7}À qual nossas doze tribos, servindo continuamente
{\ADD{a Deus}} de dia e de noite, esperam chegar; pela qual esperança, rei Agripa, eu sou acusado pelos judeus.
\VS{8}Por que se julga como incrível entre vós que Deus ressuscite aos mortos?
\VS{9}Eu realmente tinha pensado comigo mesmo, que contra o nome de Jesus eu tinha que fazer muitas oposições.
\VS{10}O que eu também fiz em Jerusalém; e tendo recebido autoridade dos chefes dos sacerdotes, eu pus em prisões a muitos dos santos; e quando eles eram mortos, eu
{\ADD{também}} dava meu voto contra
{\ADD{eles}} .
\VS{11}E tendo lhes dado punição muitas vezes por todas as sinagogas, eu os forcei a blasfemarem. E estando extremamente enfurecido contra eles, até nas cidades estrangeiras eu os persegui;
\VS{12}Nas quais, indo eu até Damasco, com autoridade e comissão dos chefes dos sacerdotes;
\VS{13}Ao meio dia, vi no caminho, rei, uma luz do céu, que brilhava muito mais que o sol, e que encheu de claridade ao redor de mim e dos que iam comigo.
\VS{14}E todos nós, tendo caído ao chão, eu ouvi uma voz que falava a mim, e dizia em língua hebraica: Saulo, Saulo, por que me persegues? Duro é para ti dar coices contra os aguilhões.
\VS{15}E eu disse: Quem és, Senhor? E ele disse: Eu sou Jesus, a quem tu persegues.
\VS{16}Mas levanta-te, e fica de pé, porque para isto eu apareci a ti, para te predeterminar como trabalhador e testemunha, tanto das coisas que
{\ADD{já}} tens visto, como das coisas que eu
{\ADD{ainda}} aparecerei a ti;
\VS{17}Livrando-te d
{\ADD{este}} povo, e
{\ADD{dos}} gentios, aos quais agora eu te envio.
\VS{18}Para abrir os olhos deles, e das trevas converterem à luz, e do poder de Satanás
{\ADD{converterem}} a Deus; para que recebam perdão dos pecados, e herança entre os santificados pela fé em mim.
\VS{19}Portanto, rei Agripa, eu não fui desobediente à visão celestial.
\VS{20}Mas sim, primeiramente aos
{\ADD{que estavam}} em Damasco e Jerusalém, e por toda a terra da Judeia, e aos gentios anunciei que se arrependessem, e se convertessem a Deus, fazendo obras dignas de arrependimento.
\VS{21}Por causa disto os judeus me pegaram no Templo, e procuravam
{\ADD{me}} matar.
\VS{22}Porém tendo eu obtido socorro de Deus, permaneço até o dia de hoje, dando testemunho tanto a pequenos como a grandes; não dizendo nada além dos que as
{\ADD{coisas}} que os profetas e Moisés tinham dito que estavam para acontecer;
\VS{23}{\ADD{Isto é}} , que o Cristo sofreria, e sendo o primeiro da ressurreição dos mortos, ia anunciar a luz a este povo e aos gentios.
\VS{24}E tendo
{\ADD{Paulo}} dito isto em
{\ADD{sua}} defesa, Festo disse em alta voz: Tu estás louco, Paulo; as muitas escrituras te fizeram enlouquecer!
\VS{25}Mas ele
{\ADD{respondeu}} : Eu não estou louco, excelentíssimo Festo; mas eu declaro palavras de verdade e de um são juízo;
\VS{26}Porque o rei, a quem eu estou falando livremente, ele sabe
{\ADD{muito bem}} destas coisas; porque eu não creio que nenhuma disto lhe seja oculto; por que isto não foi feito num canto.
\VS{27}Rei Agripa, tu crês nos profetas? Eu sei que tu crês.
\VS{28}E Agripa disse a Paulo: Por pouco tu me convences a me tornar cristão.
\VS{29}E Paulo disse: Meu desejo a Deus é que, por pouco ou por muito, não somente tu, mas todos os que estão me ouvindo hoje, tais vos tornásseis assim como eu sou, a não ser por estas correntes.
\VS{30}E tendo ele dito isto, o rei se levantou, e
{\ADD{também}} o governador, Berenice, e os que estavam sentados com eles.
\VS{31}E reunindo-se à parte, falavam uns aos outros, dizendo: Este homem nada faz
{\ADD{que seja}} digno de morte ou de prisões.
\VS{32}E Agripa disse a Festo: Este homem podia ser solto, se ele não tivesse apelado a César.

\par }\Chap{27}{\PP \VerseOne{1}E quando foi determinado que tínhamos que navegar para a Itália, entregaram a Paulo e alguns outros prisioneiros, a um centurião, por nome Júlio, do esquadrão imperial.
\VS{2}E embarcando
{\ADD{-nos}} em um navio adramitino, estando a navegar pelos lugares
{\ADD{costeiros}} da Ásia, nós partimos, estando conosco Aristarco, o macedônio de Tessalônica.
\VS{3}E no
{\ADD{dia}} seguinte, chegamos a Sidom; e Júlio, tratando bem a Paulo, permitiu
{\ADD{-lhe}} que fosse aos amigos, para
{\ADD{receber}} cuidado
{\ADD{deles}} .
\VS{4}E tendo partido dali, nós fomos navegando abaixo do Chipre, porque os ventos estavam contrários.
\VS{5}E tendo passado ao longo do mar da Cilícia e Panfília, viemos a Mira em Lícia.
\VS{6}E o centurião, tendo achado ali um navio de Alexandria, que estava navegando para a Itália, nos fez embarcar nele.
\VS{7}E indo navegando lentamente já por muitos dias, chegando com dificuldade em frente a Cnido, o vento, não nos permitindo
{\ADD{continuar por ali}} ,navegamos abaixo de Creta, em frente a Salmone.
\VS{8}E tendo com dificuldade percorrido sua costa, chegamos a um certo lugar, chamado Bons Portos, perto do qual estava a cidade de Laseia.
\VS{9}E tendo passado muito tempo, e sendo a navegação já perigosa, porque também já tinha passado o jejum, Paulo
{\ADD{os}} exortava,
\VS{10}Dizendo-lhes: Homens, eu vejo que a navegação vai ser com violência e muito dano, não somente de carga e do navio, mas também de nossas vidas.
\VS{11}Porém o centurião cria mais no capitão e no dono do navio do que no que Paulo dizia.
\VS{12}E não sendo aquele porto adequado para passar o inverno, a maioria preferiu partir dali, para ver se podiam chegar a Fênix, que é um porto de Creta, voltada para o lado do vento sudoeste e noroeste, para ali passarem o inverno.
\VS{13}E ao ventar brandamente ao sul, pareceu-lhes que eles já tinham o que queriam; e levantando a vela, foram por perto da costa de Creta.
\VS{14}Mas não muito depois houve contra ela um vento violento, chamado Euroaquilão.
\VS{15}E tendo o navio sido tomado por ele, e não podendo navegar contra o vento, nós deixamos sermos levados
{\ADD{por ele}} .
\VS{16}E correndo abaixo de uma pequena ilha, chamada Clauda, com dificuldade conseguimos manter o barquinho de reserva;
\VS{17}O qual, tendo sido levado para cima, usaram de suportes
{\ADD{para}} reforçarem o navio; e temendo irem de encontro aos bancos de areia, eles baixaram as velas e
{\ADD{deixaram}} ir à deriva.
\VS{18}E sendo muito afligidos pela tempestade, no
{\ADD{dia}} seguinte jogaram a carga para fora
{\ADD{do navio}} .
\VS{19}E no terceiro
{\ADD{dia}} , com as nossas próprias mãos jogamos fora os instrumentos do navio.
\VS{20}E não aparecendo ainda o sol, nem estrelas havia muitos dias, e sendo afligidos por não pouca tempestade, desde então tínhamos perdido toda a esperança de sermos salvos.
\VS{21}E havendo muito
{\ADD{tempo}} que não havia o que comer, então Paulo, ficando de pé no meio deles, disse: Homens, vós devíeis ter dado atenção a mim, e não terdes partido de Creta, e
{\ADD{assim}} evitar esta situação ruim e prejuízo.
\VS{22}Mas agora eu vos exorto a terdes bom ânimo; porque haverá nenhuma perda de vida de vós, além
{\ADD{somente da perda}} do navio.
\VS{23}Porque esta mesma noite esteve comigo um anjo de Deus, a quem eu pertenço e a quem eu sirvo;
\VS{24}Dizendo: Não temas, Paulo; é necessário que tu sejas apresentado a César; e eis que Deus tem te dado
{\ADD{a vida}} a todos quantos navegam contigo.
\VS{25}Portanto, homens, tende bom ânimo; porque eu creio em Deus que assim será, conforme o que me foi dito.
\VS{26}Mas é necessário que sejamos lançados a uma ilha.
\VS{27}E quando veio a décima quarta noite, sendo lançados de um lado para o outro no
{\ADD{mar}} Adriático, por volta da meia noite os marinheiros suspeitaram de que estavam se aproximando de alguma terra
{\ADD{firme}} .
\VS{28}E tendo lançado o prumo, acharam vinte braças; e passando um pouco mais adiante, voltando a lançar o prumo, acharam quinze braças.
\VS{29}E temendo de irem de encontro a lugares rochosos, lançaram da popa quatro âncoras, desejando que o dia viesse
{\ADD{logo}} .
\VS{30}E
{\ADD{quando}} os marinheiros estavam procurando fugir do navio, e baixando o barquinho de reserva ao mar, como que queriam largar as âncoras da proa,
\VS{31}Paulo disse ao centurião e aos soldados: Se estes não ficarem no navio, vós não podeis vos salvar.
\VS{32}Então os soldados cortaram os cabos do barquinho de reserva, e o deixaram cair.
\VS{33}E até enquanto o dia estava vindo, Paulo exortava a todos que comessem alguma coisa, dizendo: Hoje já é o décimo quarto dia, em que estais esperando, continuando sem comer, nada tendo experimentado.
\VS{34}Portanto eu vos exorto para que comais alguma coisa, pois é bom para vossa saúde; porque nenhum cabelo cairá de vossa cabeça.
\VS{35}E tendo dito isto, e tomando o pão, ele agradeceu a Deus na presença de todos; e partindo
{\ADD{-o}} , começou a comer.
\VS{36}E todos, tendo ficado mais encorajados, também pegaram
{\ADD{algo}} para comer.
\VS{37}E éramos todos no navio duzentas e setenta e seis almas.
\VS{38}E estando saciados de comer, eles tiraram peso do navio, lançando o trigo ao mar.
\VS{39}E tendo vindo o dia, não reconheciam a terra; mas enxergaram uma enseada que tinha praia, na qual planejaram, se pudessem, levar o navio.
\VS{40}E tendo levantado as âncoras, deixaram
{\ADD{-no}} ir ao mar, soltando também as amarras dos lemes, e levantando a vela maior ao vento, foram de levando
{\ADD{-o}} à praia.
\VS{41}Mas tendo caído em um lugar onde dois mares se encontram, encalharam ali o navio; e fixa a proa, ficou imóvel, mas a popa estava se destruindo com a força das ondas.
\VS{42}Então o conselho dos soldados foi de que matassem aos presos, para que nenhum
{\ADD{deles}} fugisse a nado.
\VS{43}Mas o centurião, querendo salvar a Paulo, impediu a intenção deles; e mandou que aqueles que pudessem nadar fossem os primeiros a se lançassem
{\ADD{ao mar}} e chegassem à terra
{\ADD{firme}} .
\VS{44}E
{\ADD{depois}} os demais, uns em tábuas, e outros em pedaços do navio. E assim aconteceu, que todos se salvaram em terra.

\par }\Chap{28}{\PP \VerseOne{1}E tendo sobrevivido, então souberam que a ilha se chamava Malta.
\VS{2}E os nativos demonstraram para conosco uma benevolência incomum; porque, tendo acendido uma fogueira, recolheram a nós todos, por causa da chuva que estava caindo, e por causa do frio.
\VS{3}E tendo Paulo recolhido uma quantidade de gravetos, e pondo-os no fogo, saiu uma víbora do calor, e fixou
{\ADD{os dentes}} na mão dele.
\VS{4}E quando os nativos vieram o animal pendurado na mão dele, disseram uns aos outros: Certamente este homem é assassino, ao qual, tendo sobrevivido do mar, a justiça não o deixa viver.
\VS{5}Porém, ele, tendo sacudido o animal ao fogo, sofreu nenhum mal.
\VS{6}E eles esperavam que ele fosse inchar, ou cair morto de repente. Mas tendo esperado muito, e vendo que nenhum incômodo tinha lhe sobrevindo, mudaram
{\ADD{de opinião}} , e diziam que ele era um deus.
\VS{7}E perto daquele mesmo lugar o homem mais importante da ilha, por nome Públio, tinha algumas propriedades; o qual nos recebeu e nos hospedou por três dias gentilmente.
\VS{8}E aconteceu que o pai de Públio estava de cama, doente de febres e disenteria; ao qual Paulo entrou, e tendo orado, pôs as mãos sobre ele, e o curou.
\VS{9}Tendo então isto acontecido, também vieram a ele outros que tinham enfermidades, e foram curados;
\VS{10}Os quais nos honraram com muitas honras; e estando
{\ADD{nós}} para navegar,
{\ADD{nos}} entregaram as coisas necessárias.
\VS{11}E três meses depois, nós partimos em um navio de Alexandria, que tinha passado o inverno na ilha; o qual tinha como símbolo os gêmeos Castor e Pólux.
\VS{12}E chegando a Siracusa, ficamos
{\ADD{ali}} três dias.
\VS{13}De onde, tendo indo ao redor da costa, chegamos a Régio; e um dia depois, ventando ao sul, viemos o segundo dia a Putéoli.
\VS{14}Onde, tendo achado
{\ADD{alguns}} irmãos, eles nos rogaram que ficássemos com eles por sete dias; e assim viemos a Roma.
\VS{15}E os irmãos, ao ouvirem
{\ADD{notícias}} sobre nós, desde lá nos saíram ao encontro até a praça de Ápio, e as três tavernas; e Paulo, tendo os visto, agradeceu a Deus, e tomou coragem.
\VS{16}E quando chegamos a Roma, o centurião entregou os prisioneiros ao chefe da guarda, mas a Paulo foi permitido morar por si mesmo à parte, junto com o soldado que o guardava.
\VS{17}E aconteceu que, três dias depois, Paulo chamou juntos os chefes dos judeus; e ao se reunirem, disse-lhes: Homens irmãos, tendo eu nada feito contra o povo, ou contra os costumes dos pais,
{\ADD{mesmo assim}} eu vim preso desde Jerusalém, entregue em mãos dos romanos.
\VS{18}Os quais, tendo me investigado, queriam
{\ADD{me}} soltar, por não haver em mim nenhum crime de morte.
\VS{19}Mas os judeus, dizendo em contrário, eu fui forçado a apelar a César;
{\ADD{mas}} não como que eu tenha que acusar a minha nação.
\VS{20}Então por esta causa eu vos chamei até mim, para
{\ADD{vos}} ver e falar; porque pela esperança de Israel eu estou agora preso nesta corrente.
\VS{21}Mas eles lhe disseram: Nós nem recebemos cartas da Judeia relacionadas a ti, nem algum dos irmãos, tendo vindo aqui, tem nos informado ou falado de ti algum mal.
\VS{22}Mas nós queríamos ouvir de ti o que tu pensas; porque, quanto a esta seita, conhecemos que em todo lugar
{\ADD{há quem}} fale contra ela.
\VS{23}E tendo eles lhe determinado um dia, muitos vieram até onde ele estava morando; aos quais ele declarava e dava testemunho do Reino de Deus; e procurava persuadi-los quanto a Jesus, tanto pela Lei de Moisés, como
{\ADD{pelos}} profetas, desde a manhã até a tarde.
\VS{24}E alguns criam nas coisas que ele dizia; mas outros não criam.
\VS{25}E estando discordantes entre si, despediram-se, tendo Paulo disto
{\ADD{esta}} palavra: O Espírito Santo corretamente falou a nossos pais por meio de Isaías o profeta,
\VS{26}Dizendo: Vai a este povo, e dize: De fato ouvireis, mas de maneira nenhuma entendereis; e de fato vereis, mas de maneira nenhuma enxergareis.
\VS{27}Porque o coração deste povo está insensível, seus ouvidos ouvem com dificuldade, e seus olhos estão fechados; para que em maneira nenhuma vejam com os olhos, nem ouçam com os ouvidos, nem entendam com o coração, e se convertam, e eu os cure.
\VS{28}Portanto seja conhecido por vós que a salvação de Deus foi enviada aos gentios; e eles
{\ADD{a}} ouvirão.
\VS{29}E havendo ele dito isto, os judeus foram embora, havendo entre eles grande discussão.
\VS{30}E Paulo ficou dois anos inteiros em sua própria casa alugada; e recebia a todos quantos vinham a ele;
\VS{31}Pregando o Reino de Deus, e ensinando com ousadia a doutrina do Senhor Jesus Cristo, sem impedimento algum.\par }